\documentclass[sigconf,nonacm]{acmart}

\settopmatter{printfolios=true}

% ---------- Packages ----------
\usepackage{microtype}
\usepackage{enumitem}

% ---------- Dense formatting ----------
\setlist{nosep,leftmargin=*}
% Tighten space around theorem environments
\makeatletter
\def\thm@space@setup{%
  \thm@preskip=4pt plus 1pt minus 1pt
  \thm@postskip=4pt plus 1pt minus 1pt
}
\makeatother
\hypersetup{hidelinks}

% ---------- Diagram support ----------
\usepackage{tikz}
\usetikzlibrary{positioning,calc,arrows.meta,fit}

% ---------- Theorem environments ----------
\theoremstyle{definition}
\newtheorem{definition}{Definition}[section]
\newtheorem{example}[definition]{Example}
\newtheorem{question}[definition]{Question}

\theoremstyle{plain}
\newtheorem{theorem}[definition]{Theorem}
\newtheorem{lemma}[definition]{Lemma}
\newtheorem{proposition}[definition]{Proposition}
\newtheorem{corollary}[definition]{Corollary}
\newtheorem{conjecture}[definition]{Conjecture}
\newtheorem{openproblem}[definition]{Open Problem}

\theoremstyle{remark}
\newtheorem{remark}[definition]{Remark}

% ---------- Internal-draft helpers ----------
\newcommand{\TODO}[1]{\medskip\noindent\fbox{\parbox{\dimexpr\linewidth-2\fboxsep}{\textbf{TODO:}\ #1}}\medskip}

\title{Semantic Schemas and Bellman's Principle}

\author{Joseph M. Hellerstein}
\affiliation{%
  \institution{UC Berkeley \& Amazon Web Services}
  \city{Berkeley, CA}
  \country{USA}
}

\author{Sarah Morin}
\affiliation{%
  \institution{UC Berkeley}
  \city{Berkeley, CA}
  \country{USA}
}

\author{Max Willsey}
\affiliation{%
  \institution{UC Berkeley}
  \city{Berkeley, CA}
  \country{USA}
}

\begin{document}

\maketitle

%%%%%%%%%%%%%%%%%%%%%%%%%%%%%%%%%%%%%%%%%%%%%%%%%%%%%%%%%%%%%%%
\section{Introduction}
\label{sec:introduction}
%%%%%%%%%%%%%%%%%%%%%%%%%%%%%%%%%%%%%%%%%%%%%%%%%%%%%%%%%%%%%%%

Many optimization algorithms exploit compositional structure.

A query optimizer decomposes a query into subqueries.
A dynamic program decomposes a problem into subproblems.
An equality-saturation engine decomposes an expression into equivalent
subexpressions.

Although these systems appear quite different, they all confront the
same question:

\begin{quote}
\emph{What information about a solved subproblem must be retained so
that larger problems can still be solved correctly?}
\end{quote}

Retaining every distinction among subproblem realizations is always
correct but quickly becomes infeasible.
Discarding distinctions too aggressively may destroy globally optimal
solutions.

The success of a compositional optimizer therefore depends on choosing
the right semantic interface between parent and child subproblems.

Relational query optimization provides the classical example.
For a fixed relational expression, many physical plans compute the same
result.
Some nevertheless differ in properties, such as output order, that may
be exploited by enclosing operators.
Selinger's notion of \emph{interesting orders} preserves exactly those
distinctions that remain observable during future composition.

Equality saturation faces the same issue from the opposite direction.
An e-class intentionally forgets distinctions among equivalent
expressions.
Extraction then assumes that one representative of each e-class
suffices.
When contextual semantic distinctions matter, however, a single
equivalence class may be too coarse.

These observations suggest that the central object of compositional
optimization is neither a search algorithm nor an objective function.
It is the semantic interface through which subproblems compose.

This paper develops a mathematical framework for those interfaces.

The framework begins with a concrete semantic schema describing all
realizations of a compositional problem.
Semantic quotients identify realizations that are interchangeable at
different resolutions.
These quotients are organized into hierarchies of
\emph{multe-classes}.
Non-atomic realizations reference child multe-classes rather than
individual child realizations, producing a \emph{multe-graph}.

A problem instance is obtained by imposing requirements on this
semantic structure.
Instantiation propagates those requirements through the multe-graph,
removing alternatives and semantic distinctions that become
unobservable.
Canonicalization then chooses one representative for each surviving
multe-class.

This separation isolates semantic representation from optimization.

Optimization supplies only a preference among representatives.
Scalar optimization, Pareto optimization, and other optimization
criteria therefore become canonicalization policies over the same
underlying semantic interfaces.

Within this framework, Bellman's Principle admits a simple semantic
interpretation.

A parent interacts with each child only through the semantic interface
represented by the referenced multe-class.
If that interface is sufficient for future composition, and the
optimization objective is monotone with respect to it, then one optimal
representative per multe-class suffices for global optimization.

Bellman's Principle is therefore not an independent assumption.
It is a theorem about semantic interfaces.

This perspective also separates optimization from search.

Branch-and-bound, for example, computes monotone entailments about
unknown optimum values while leaving the semantic representation
unchanged.
Likewise, changing the optimization objective leaves the semantic
interfaces unchanged.
The framework therefore cleanly separates semantic representation,
problem instantiation, optimization, canonicalization, and search.

This separation is not merely aesthetic; its dividends organize the
paper.

First, correctness becomes modular.  Each layer carries its own
obligation---representation must preserve observable behavior,
instantiation must preserve satisfying realizations, optimization
must compose across interfaces, search must preserve the
optimum---and each can be discharged, or generalized, without
disturbing the others.  Additive cost relaxes to Pareto frontiers
without touching the semantic representation; acyclic schemas relax
to fixed points without touching the objective.

Second, systems that grew up in different communities become
configurations of one theory.  Selinger's optimizer, Cascades,
e-graph extraction, and Knuth--Bendix completion differ in which
layers they fix and which they vary; stated over shared interfaces,
results transfer among them.

Third, and most consequentially, the separation converts design art
into mathematics.  Only after the semantic interface is isolated
from the objective and the search strategy can one ask where
interfaces come from and what they cost.  The payoffs of
Section~\ref{sec:payoffs} answer both: interfaces can be derived
rather than designed, and their unavoidable size is an intrinsic
complexity measure of the problem, not an artifact of any
algorithm.

The remainder of the paper develops this framework in a
simplest-first manner.
We begin with finite acyclic semantic schemas, tree-shaped refinement
hierarchies, Boolean local requirements, and additive cost.
The later sections show that these assumptions may each be relaxed
independently.

%%%%%%%%%%%%%%%%%%%%%%%%%%%%%%%%%%%%%%%%%%%%%%%%%%%%%%%%%%%%%%%
\subsection{Roadmap}
\label{sec:intro-roadmap}
%%%%%%%%%%%%%%%%%%%%%%%%%%%%%%%%%%%%%%%%%%%%%%%%%%%%%%%%%%%%%%%

Section~\ref{sec:defining-multegraph} defines the multe-graph
carefully and from scratch: semantic schemas
(Section~\ref{chap:schemas}), semantic quotients and their
refinement hierarchies (Section~\ref{chap:quotients}), and the
combination of the two (Section~\ref{sec:multe-graphs}).  A reading
guide at its head summarizes the departures from an ordinary e-graph
and makes the section skimmable for readers who already know
e-graphs.
Section~\ref{sec:instances} defines requirements and problem
instances.
Section~\ref{sec:canonicalization} separates canonicalization
from instantiation.
Section~\ref{sec:extraction} defines extraction as requirement
propagation followed by canonicalization.
The middle sections develop examples, theorem statements, and
generalizations.
Section~\ref{sec:payoffs} collects the intended headline results of
the outbound article: the derivation of semantic interfaces from
continuation behavior, the width and dimension of a compositional
optimization problem, and applications to search and to
rewriting-system completion.

%%%%%%%%%%%%%%%%%%%%%%%%%%%%%%%%%%%%%%%%%%%%%%%%%%%%%%%%%%%%%%%
\section{The Multe-Graph, Carefully}
\label{sec:defining-multegraph}
%%%%%%%%%%%%%%%%%%%%%%%%%%%%%%%%%%%%%%%%%%%%%%%%%%%%%%%%%%%%%%%

This section defines the multe-graph slowly and from scratch,
because the outbound article must serve readers who know neither
e-graphs nor query optimizers.  Readers of this internal draft know
both, so the section is built to be skimmed: everything downstream
depends only on its definitions and notation, never on the
surrounding scaffolding and examples.

\paragraph{Reading guide.}
For a reader who already knows e-graphs, the genuine departures are
four.

\begin{enumerate}

\item
Composition is presented as a relation \(\mathcal R_d\) over
e-nodes, deliberately stripped of evaluation direction and
preference (Section~\ref{sec:composition}).

\item
A goal carries not one equivalence but a \emph{refinement hierarchy}
of semantic equivalences; multe-classes are the blocks of these
quotients, and a \emph{cut} assembles a partition from classes at
different resolutions in different branches
(Section~\ref{sec:refinement-hierarchies}).

\item
A non-atomic e-node references one multe-class per child slot, and
the composition relation is required to agree with those references
exactly: the e-node appears in the relation once for every
combination of members drawn from its referenced classes, and for no
other child tuples.  (An e-node whose admissible children lack this
shape is split into several e-nodes that have it.)  The consequence
is that slot-wise substitution always works: any member of a
referenced class may replace any other, without consulting the
remaining slots.  It also makes composition functional over the
quotient: \(\chi\) assigns each e-node its tuple of referenced
classes (Section~\ref{sec:multe-graph-definition}).

\item
The notation to retain: a schema
\(\mathcal S=(X,\mathsf N,\Sigma,\mathcal R)\), a multe-graph
\(\mathcal G=(\mathcal S,\mathcal C,\chi)\), the typing judgment
\(n:x\), and multe-classes \([n]_{\equiv}\).

\end{enumerate}

An ordinary e-graph is the special case in which every refinement
hierarchy is trivial: one universal class per goal.  If the four
points above read as expected, the rest of this section can be
treated as reference material.

%%%%%%%%%%%%%%%%%%%%%%%%%%%%%%%%%%%%%%%%%%%%%%%%%%%%%%%%%%%%%%%
\subsection{Semantic Schemas}
\label{chap:schemas}
%%%%%%%%%%%%%%%%%%%%%%%%%%%%%%%%%%%%%%%%%%%%%%%%%%%%%%%%%%%%%%%
A \emph{semantic schema} is the object at the base of the framework:
a description of all admissible realizations of a compositional
problem, prior to any particular problem instance, cost model, or
search strategy.  The name is chosen by analogy with database
schemas.  A schema describes a family of related problems; an
individual problem instance is obtained later, by registering
requirements against the schema (Section~\ref{sec:instances}).

The formal definition (Definition~\ref{def:concrete-schema}) arrives
at the end of this subsection, once its ingredients are in place.
We build up to it in stages, beginning with the unconstrained local
vocabulary of a schema: goals and their e-nodes.  Semantic
equivalence and compositional references will be added in the
subsections that follow.  No requirements, costs, or
canonicalization are assumed.

%%%%%%%%%%%%%%%%%%%%%%%%%%%%%%%%%%%%%%%%%%%%%%%%%%%%%%%%%%%%%%%
\subsubsection{Goals and E-Nodes}
\label{sec:goals-enodes}
%%%%%%%%%%%%%%%%%%%%%%%%%%%%%%%%%%%%%%%%%%%%%%%%%%%%%%%%%%%%%%%

We begin with the local objects from which a semantic schema is
assembled.

At this stage, we deliberately ignore equivalence, composition,
requirements, costs, and canonicalization.  We assume only that a
problem can be divided into named goals, and that each goal admits a
collection of alternative realizations.

\begin{definition}[Goal]
A \emph{goal} is a semantic task whose admissible realizations are to
be represented by the schema.
\end{definition}

The definition is deliberately thin, so it helps to say plainly what
work the word does.  A goal poses a problem or subproblem---\emph{what}
to be realized---without
solving it; the e-nodes defined below are the candidate solutions---\emph{how}
to achieve the goal---that are on offer for it.  In dynamic-programming terms, a goal is a key of the
memo table and its e-nodes are the entries stored under that key; in
typed terms, a goal is a type and its e-nodes are the inhabitants.
In every register, the goal names the question and the e-nodes are
its answers.

Depending on the
application, a goal may be

\begin{itemize}
\item a relational expression to be implemented;
\item a program expression to be rewritten;
\item a logical proposition to be established;
\item an output value to be produced; or
\item an abstract capability to be provided.
\end{itemize}

Let \(X\) be a finite set of goals.  For each \(x \in X\), let
\(\mathsf{N}(x)\) denote the finite set of e-nodes associated with
\(x\).

\begin{definition}[E-node]
An \emph{e-node} for a goal \(x\) is one admissible local realization
of \(x\).
\end{definition}

Thus every e-node belongs to exactly one goal.  We write \(n : x\)
to mean that \(n \in \mathsf{N}(x)\).

The notation anticipates a reading we develop at the end of the
document (Section~\ref{sec:types}): goals are base types, e-nodes
are their constructors, and \(n:x\) is a typing judgment.

An e-node is a local realization of its goal.  A non-atomic e-node
specifies its operator and the semantic requirements on its children,
without selecting particular child realizations.

We distinguish two cases.

\begin{definition}[Atomic and non-atomic e-nodes]
An e-node is \emph{atomic} if its realization requires no subordinate
goals.  It is \emph{non-atomic} otherwise.
\end{definition}

This distinction records whether an alternative is compositionally
closed, but does not yet specify its children.

\begin{example}[Relational query optimization]
Consider the goal \(x_{R \Join S}\) of producing the relation
\(R \Join S\).  Possible e-nodes for this goal include
\(\mathsf{HashJoin}\), \(\mathsf{MergeJoin}\), and
\(\mathsf{NestedLoopJoin}\).

These are different local realizations of the same goal.  Each is
non-atomic: realizing the join requires subordinate realizations of
the goals \(x_R\) and \(x_S\).  The precise relationship between the
join e-node and those subordinate goals will be introduced in
Section~\ref{sec:composition}.

A table scan for \(R\), by contrast, may be represented by an atomic
e-node for the goal \(x_R\).
\end{example}

\begin{example}[Expression optimization]
Let \(x_e\) be the goal of realizing an expression semantically
equivalent to \(a(b+c)\).  Possible non-atomic e-nodes for this goal
include \(\mathsf{Mul}\) and \(\mathsf{Add}\).  The first represents
a multiplicative decomposition of the expression; the second
represents the distributed additive decomposition \(ab+ac\).

At this stage we record only that these are distinct local alternatives
for the same goal.  Their subordinate expressions, and the semantic
classes to which those subordinate expressions belong, will be
introduced with composition and multe-graphs.
\end{example}

At this point, the schema contains only the disjoint family
\(\{\mathsf{N}(x)\}_{x\in X}\) of alternatives associated with each
goal.

Nothing yet identifies two e-nodes as equivalent, and nothing yet says
how a non-atomic e-node depends on subordinate alternatives.  Those two
pieces of structure will be introduced separately:

\begin{enumerate}
\item multe-classes will identify semantically interchangeable
      alternatives; and
\item the multe-graph will allow non-atomic e-nodes to reference
      multe-classes of subordinate goals.
\end{enumerate}

Keeping these steps separate is important.  Goals and e-nodes describe
the local vocabulary of the schema.  Equivalence and composition are
additional structures imposed on that vocabulary.

%%%%%%%%%%%%%%%%%%%%%%%%%%%%%%%%%%%%%%%%%%%%%%%%%%%%%%%%%%%%%%%
\subsubsection{Composition}
\label{sec:composition}
%%%%%%%%%%%%%%%%%%%%%%%%%%%%%%%%%%%%%%%%%%%%%%%%%%%%%%%%%%%%%%%

Goals and e-nodes describe the local alternatives available to a
schema.  We next describe how alternatives for one goal may be composed
from alternatives for other goals.

We begin with a concrete, unquotiented account of composition.  This
will later provide the semantic relation represented more compactly by
references to multe-classes.

Let \(X\) be the set of goals, and let \(\mathsf{N}(x)\) be the set of
e-nodes for each goal \(x \in X\).

\begin{definition}[Composition signature]
A \emph{composition signature} is a tuple
\[
d=(x \Rightarrow x_1,\ldots,x_k)
\]
consisting of a parent goal \(x\) and an ordered list of subgoals
\(x_1,\ldots,x_k\).
\end{definition}

A composition signature specifies only the types of the participating
alternatives.  It does not yet say which combinations are semantically
admissible.

\begin{definition}[Composition relation]
For a composition signature \(d=(x \Rightarrow x_1,\ldots,x_k)\),
a \emph{composition relation} is a relation
\[
\mathcal{R}_d
\subseteq
\mathsf{N}(x)
\times
\mathsf{N}(x_1)
\times\cdots\times
\mathsf{N}(x_k).
\]

A tuple \((n,n_1,\ldots,n_k)\in\mathcal{R}_d\) states that the
parent alternative \(n\) may be realized using the subgoal
alternatives \(n_1,\ldots,n_k\).
\end{definition}

Composition is relational rather than functional. That is,
for a given parent goal \(x\), there may be many admissible
configurations of subgoals.  These may differ in both the parent
alternative and the alternatives chosen for its subgoals:
\[
(n,n_1,\ldots,n_k)\in\mathcal{R}_d,
\qquad
(n',n'_1,\ldots,n'_k)\in\mathcal{R}_{d'}.
\]

For example, the goal \(R\Join S\) may be realized by a hash join over
arbitrary implementations of \(R\) and \(S\), or by a merge join over
implementations that provide suitable orders.

Conversely, one parent alternative may admit several child
realizations:
\[
(n,n_1,\ldots,n_k) \in \mathcal{R}_d,
\qquad
(n,n'_1,\ldots,n'_k) \in \mathcal{R}_d.
\]
The parent e-node records a local semantic alternative, not necessarily
a unique derivation from concrete children.

In most presentations these ingredients arrive fused.  A
dynamic-programming recurrence defines the admissible compositions,
the order in which they are evaluated, and the minimization over
them, all in a single equation.  An oriented rewrite rule asserts an
equation, a direction of travel, and a preference for one side, all
in one arrow.  An optimizer's transformation rules are typically
written to fire in whatever direction its search happens to proceed.
The relational treatment separates three matters that these
artifacts conflate:

\begin{enumerate}

\item
\emph{semantic admissibility}: whether a tuple belongs to
\(\mathcal{R}_d\);

\item
\emph{evaluation direction}: whether an implementation explores the
relation from parent to children, from children to parent, or in both
directions; and

\item
\emph{preference}: whether one admissible tuple is cheaper or otherwise
preferred to another.

\end{enumerate}

Only the first belongs to the semantic schema.  A composition relation
has no intrinsic traversal direction and contains no cost information.
The other two return later as interchangeable policies over the same
schema: evaluation direction as a search strategy---bottom-up
Selinger versus top-down Cascades (Section~\ref{sec:demand})---and
preference as a canonicalization policy
(Section~\ref{sec:canonicalization}).  Keeping them out of the
schema now is what makes them interchangeable later.

\begin{example}[Join composition]
Let \(x_R\), \(x_S\), and \(x_{R \Join S}\) be goals, and consider
the composition signature
\(d_{\mathsf{join}} = ( x_{R \Join S} \Rightarrow x_R,x_S )\).
The relation
\[
\mathcal{R}_{d_{\mathsf{join}}}
\subseteq
\mathsf{N}(x_{R \Join S})
\times
\mathsf{N}(x_R)
\times
\mathsf{N}(x_S)
\]
may contain tuples such as \((\mathsf{HashJoin},r,s)\),
\((\mathsf{MergeJoin},r',s')\), and
\((\mathsf{NestedLoopJoin},r'',s'')\).

Membership states only that the corresponding child alternatives may
be combined to realize the parent alternative.  It does not rank the
join methods or prescribe an order in which they are considered.
\end{example}

\begin{example}[Expression composition]
For a binary addition signature \(d_{+} = (x \Rightarrow x_1,x_2)\),
the relation
\(\mathcal{R}_{d_{+}}
\subseteq
\mathsf{N}(x)\times\mathsf{N}(x_1)\times\mathsf{N}(x_2)\)
records which child alternatives may be combined by addition to realize
an alternative for \(x\).

Equivalent parent expressions may arise from different child tuples,
and the same child tuple may participate in several equivalent parent
representations.
\end{example}

\begin{definition}[Concrete semantic schema]
\label{def:concrete-schema}
A \emph{concrete semantic schema} is a tuple
\[
\mathcal S=(X,\mathsf N,\Sigma,\mathcal R)
\]
consisting of

\begin{enumerate}

\item
a finite set of goals \(X\);

\item
a finite set of e-nodes \(\mathsf{N}(x)\) for each \(x \in X\);

\item
a finite collection \(\Sigma\) of composition signatures; and

\item
a composition relation \(\mathcal{R}_d\) for each signature
\(d\in\Sigma\).

\end{enumerate}
\end{definition}

The concrete schema may be much larger than an optimizer should
materialize.  It distinguishes individual child e-nodes even when a
parent can observe only some coarser semantic property of them.

The next subsection introduces multe-classes.  Once those classes
are available, a non-atomic parent e-node will refer to child
multe-classes rather than to individual child e-nodes.  The quotient
improves the form of composition as well as its size: the relation
\(\mathcal{R}_d\), which ranges over tuples of individual e-nodes,
becomes a \emph{function} assigning each parent e-node one tuple of
child multe-classes (splitting e-nodes where necessary so that each
admits a single tuple).  Composition is relational over e-nodes but
functional over the right quotient.  The resulting multe-graph is a
quotient representation of the concrete composition relations
defined here.

%%%%%%%%%%%%%%%%%%%%%%%%%%%%%%%%%%%%%%%%%%%%%%%%%%%%%%%%%%%%%%%
\subsection{Semantic Quotients}
\label{chap:quotients}
%%%%%%%%%%%%%%%%%%%%%%%%%%%%%%%%%%%%%%%%%%%%%%%%%%%%%%%%%%%%%%%

%%%%%%%%%%%%%%%%%%%%%%%%%%%%%%%%%%%%%%%%%%%%%%%%%%%%%%%%%%%%%%%
\subsubsection{Multe-Classes}
\label{sec:multe-classes}
%%%%%%%%%%%%%%%%%%%%%%%%%%%%%%%%%%%%%%%%%%%%%%%%%%%%%%%%%%%%%%%

The concrete semantic schema distinguishes every e-node associated
with a goal.  An optimizer generally need not preserve all of these
distinctions.  It may instead treat e-nodes as interchangeable whenever
the surrounding computation cannot observe their differences.

We first consider a single semantic equivalence for a single goal.
Families of equivalences ordered by refinement will be introduced in
Section~\ref{sec:refinement-hierarchies}.

Let \(x\in X\) be a goal, with e-node set \(\mathsf{N}(x)\).

\begin{definition}[Semantic equivalence]
A \emph{semantic equivalence} for a goal \(x\) is an equivalence
relation
\[
{\equiv}\ \subseteq\ \mathsf{N}(x)\times\mathsf{N}(x).
\]

For \(n,n'\in\mathsf{N}(x)\), the judgment \(n\equiv n'\) means that
the distinction between \(n\) and \(n'\) is not retained by this
equivalence.
\end{definition}

\begin{definition}[Multe-class]
Given a semantic equivalence \(\equiv\) on \(\mathsf{N}(x)\), the
\emph{multe-class} of \(n\in\mathsf{N}(x)\) is
\[
[n]_{\equiv}
=
\{n'\in\mathsf{N}(x)\mid n'\equiv n\}.
\]
\end{definition}

The quotient set \(\mathsf{N}(x)/{\equiv}\) is the collection of
multe-classes induced by \(\equiv\).

We refer to the choice of equivalence, informally, as a
\emph{resolution}, by analogy with images: a finer equivalence
retains more distinctions, a coarser one fewer.

A multe-class is therefore not an arbitrary collection of e-nodes.
It is one block of a semantic quotient of the alternatives for a
single goal.

\begin{example}[Output order]
Let \(x_R\) be the goal of producing relation \(R\).
Suppose the semantic equivalence \(\equiv_{\mathsf{order}}\)
distinguishes e-nodes according to the output order they provide:
two implementations satisfy \(n\equiv_{\mathsf{order}}n'\) when they
provide the same relevant output order.

The quotient may contain classes such as
\[
[R]_{\mathsf{sorted}(a)},
\qquad
[R]_{\mathsf{sorted}(b)},
\qquad
[R]_{\mathsf{unordered}}.
\]

Each class may contain many physical implementations.
For example, two index scans using different indexes may belong to the
same multe-class if both provide order on \(a\) and no distinction
retained by \(\equiv_{\mathsf{order}}\) separates them.
\end{example}

\begin{example}[Coarser logical equivalence]
At a coarser resolution \(\equiv_{\mathsf{logical}}\), output order
may be ignored entirely.  Then all implementations that produce
relation \(R\) belong to one class:
\([R]_{\mathsf{any}} = \mathsf{N}(x_R)\).
\end{example}

The distinction between a goal and a multe-class is important.

A goal identifies the task to be realized.
A multe-class identifies a set of realizations that are
interchangeable at a particular semantic resolution.
All members of \([R]_{\mathsf{sorted}(a)}\) realize the goal
\(x_R\), but they additionally agree on a property that may matter
to composition.

Likewise, the distinction between a multe-class and a preferred
representative is important.
Membership in a multe-class expresses semantic interchangeability at
one resolution.
It says nothing about cost and chooses no representative.
Selection among the members of a multe-class is deferred to
canonicalization.

The multe-classes induced by a single semantic equivalence
\(\equiv\) form a partition of \(\mathsf{N}(x)\):
\[
\mathsf{N}(x)
=
\biguplus_{C\in \mathsf{N}(x)/{\equiv}} C.
\]

A single semantic equivalence gives one partition of
\(\mathsf{N}(x)\).  Our intended structure contains several such
equivalences, ordered by refinement.  A finer equivalence preserves
distinctions that a coarser equivalence forgets.

We next define the compatibility conditions among these
equivalences and relate them to a tree-shaped representation of nested
multe-classes.  In that representation, any cut that covers all leaves
corresponds to one partition of \(\mathsf{N}(x)\); the tree as a whole
does not constitute a single partition.

%%%%%%%%%%%%%%%%%%%%%%%%%%%%%%%%%%%%%%%%%%%%%%%%%%%%%%%%%%%%%%%
\subsubsection{Refinement Hierarchies}
\label{sec:refinement-hierarchies}
%%%%%%%%%%%%%%%%%%%%%%%%%%%%%%%%%%%%%%%%%%%%%%%%%%%%%%%%%%%%%%%

A single semantic equivalence determines one partition of the e-nodes
for a goal.  Our intended structure retains several such equivalences,
ranging from fine distinctions to coarse ones.

We begin with the simplest case: a finite linear hierarchy of semantic
equivalences.  This assumption produces a tree-shaped hierarchy of
multe-classes.  More general partially ordered families of quotients
will be considered later.

Let \(x\in X\) be a goal, with e-node set \(\mathsf{N}(x)\).

\begin{definition}[Refinement]
Let \(\equiv_1\) and \(\equiv_2\) be semantic equivalences on
\(\mathsf{N}(x)\).

We say that \(\equiv_1\) \emph{refines} \(\equiv_2\), written
\(\equiv_1 \preceq \equiv_2\), when
\[
n\equiv_1 n'
\quad\Longrightarrow\quad
n\equiv_2 n'
\]
for all \(n,n'\in\mathsf{N}(x)\).  Equivalently,
\({\equiv_1}\subseteq{\equiv_2}\) as binary relations.
\end{definition}

Thus \(\equiv_1\preceq\equiv_2\) means that \(\equiv_1\) is finer:
it may distinguish e-nodes that \(\equiv_2\) identifies.

Refinement can also be stated in terms of classes.

\begin{proposition}
If \(\equiv_1\preceq\equiv_2\), then every multe-class induced by
\(\equiv_1\) is contained in a unique multe-class induced by
\(\equiv_2\).
\end{proposition}

\begin{proof}
Let \(C=[n]_{\equiv_1}\).
For every \(n'\in C\), we have \(n'\equiv_1 n\), and therefore
\(n'\equiv_2 n\).  Hence \(C\subseteq[n]_{\equiv_2}\).
Uniqueness follows because the \(\equiv_2\)-classes form a partition of
\(\mathsf{N}(x)\).
\end{proof}

This containment property allows classes belonging to different
equivalences to be organized into a hierarchy.

\begin{definition}[Linear refinement hierarchy]
A \emph{linear refinement hierarchy} for a goal \(x\) is a finite
sequence of semantic equivalences
\[
\equiv_0
\preceq
\equiv_1
\preceq
\cdots
\preceq
\equiv_h
\]
on \(\mathsf{N}(x)\), where the coarsest equivalence
\(\equiv_h\) is universal:
\[
n\equiv_h n'
\qquad
\text{for all }n,n'\in\mathsf{N}(x).
\]

The equivalence \(\equiv_0\) is the finest equivalence represented by
the hierarchy.
It need not be equality.
\end{definition}

The universal equivalence contributes the root class
\(\mathsf{N}(x)\).
Every class at a finer level is contained in a unique class at every
coarser level.

\begin{definition}[Tree of multe-classes]
Given a linear refinement hierarchy
\(\equiv_0\preceq\cdots\preceq\equiv_h\),
its \emph{multe-class tree} is the inclusion hierarchy formed by all
distinct classes appearing in the quotients
\[
\mathsf{N}(x)/{\equiv_0},
\ldots,
\mathsf{N}(x)/{\equiv_h}.
\]

The root is \(\mathsf{N}(x)\).
The parent of a non-root class \(C\) is its least strict superclass
among the represented multe-classes.
\end{definition}

Because the equivalences are linearly ordered by refinement, every
non-root class has a unique parent.  The resulting inclusion structure
is therefore a rooted tree after duplicate classes are identified.

The leaves need not be singletons.  If the finest represented
equivalence \(\equiv_0\) still identifies several e-nodes, then those
e-nodes remain together in a leaf multe-class.  The hierarchy records
only the distinctions represented by the schema, not necessarily every
distinction among concrete e-nodes.

\begin{example}[Output-order hierarchy]
Let \(x_R\) be the goal of producing relation \(R\).

At a fine equivalence \(\equiv_{\mathsf{order}}\), implementations are
identified only when they provide the same relevant output order.  The
resulting quotient may contain
\[
[R]_{\mathsf{sorted}(a)},
\qquad
[R]_{\mathsf{sorted}(b)},
\qquad
[R]_{\mathsf{unordered}}.
\]

At a coarser equivalence \(\equiv_{\mathsf{logical}}\), output order
is ignored: \([R]_{\mathsf{any}} = \mathsf{N}(x_R)\).

We then have
\(\equiv_{\mathsf{order}} \preceq \equiv_{\mathsf{logical}}\),
and correspondingly each of
\([R]_{\mathsf{sorted}(a)}\), \([R]_{\mathsf{sorted}(b)}\), and
\([R]_{\mathsf{unordered}}\) is contained in
\([R]_{\mathsf{any}}\).
\end{example}

A single e-node therefore belongs to one class at each equivalence in
the hierarchy.  For example, an implementation that produces order on
\(a\) belongs both to \([R]_{\mathsf{sorted}(a)}\) under
\(\equiv_{\mathsf{order}}\), and to \([R]_{\mathsf{any}}\) under
\(\equiv_{\mathsf{logical}}\).

These memberships are not competing classifications.  They describe
the same e-node at different degrees of semantic distinction.

\begin{definition}[Cut]
A \emph{cut} of a multe-class tree is a collection \(\mathcal{K}\)
of pairwise incomparable multe-classes such that every leaf is
contained in exactly one member of \(\mathcal{K}\).
\end{definition}

Every cut partitions the e-node set:
\[
\mathsf{N}(x)
=
\biguplus_{C\in\mathcal{K}} C.
\]

A cut need not consist of all classes from one equivalence in the
original linear hierarchy.  It may select a fine class in one branch
and a coarse class in another.  Such nonuniform cuts will become
important when requirements preserve some semantic distinctions while
allowing others to collapse.

The distinction between the hierarchy and a cut is therefore
fundamental.

The hierarchy records every represented semantic distinction.
A cut records one current partition of the alternatives obtained by
choosing which distinctions remain visible in each branch.

So far, each goal has been considered in isolation.
The next subsection combines these refinement hierarchies with
composition.  A parent e-node will refer to child multe-classes at the
coarsest points in their hierarchies sufficient for that parent
alternative.

%%%%%%%%%%%%%%%%%%%%%%%%%%%%%%%%%%%%%%%%%%%%%%%%%%%%%%%%%%%%%%%
\subsection{Multe-Graphs}
\label{sec:multe-graphs}
%%%%%%%%%%%%%%%%%%%%%%%%%%%%%%%%%%%%%%%%%%%%%%%%%%%%%%%%%%%%%%%

A concrete semantic schema describes composition in terms of
individual e-nodes.  A refinement hierarchy describes which
distinctions among those e-nodes may be forgotten.  A
\emph{multe-graph} combines these two structures.

As in an ordinary e-graph, a non-atomic e-node references classes of
possible children rather than selecting particular child e-nodes.
Unlike an ordinary e-graph, each referenced class may occur at any
represented point in its goal's refinement hierarchy.

%%%%%%%%%%%%%%%%%%%%%%%%%%%%%%%%%%%%%%%%%%%%%%%%%%%%%%%%%%%%%%%
\subsubsection{Definition}
\label{sec:multe-graph-definition}
%%%%%%%%%%%%%%%%%%%%%%%%%%%%%%%%%%%%%%%%%%%%%%%%%%%%%%%%%%%%%%%

For each goal \(x\), let \(\mathcal C(x)\) denote the collection of
all multe-classes appearing in its refinement hierarchy.

Let \(n\in\mathsf N(x)\) be a non-atomic e-node with composition
signature \(\sigma(n) = d = (x\Rightarrow x_1,\ldots,x_k)\).

Its concrete composition fiber is
\[
\mathcal R_d[n]
=
\left\{
(n_1,\ldots,n_k)
\;\middle|\;
(n,n_1,\ldots,n_k)\in\mathcal R_d
\right\}.
\]

A \emph{fiber} is simply an inverse image at a point:
\(\mathcal R_d[n]\) collects the child tuples that the relation
pairs with \(n\).

\begin{definition}[Child references]
The \emph{child references} of \(n\) are a typed tuple
\[
\chi(n)
=
(C_1,\ldots,C_k)
\in
\mathcal C(x_1)\times\cdots\times\mathcal C(x_k)
\]
such that
\[
\mathcal R_d[n]
=
C_1\times\cdots\times C_k.
\]

An atomic e-node has the empty tuple of child references.
\end{definition}

Thus each non-atomic e-node has one fixed tuple of referenced
multe-classes.  The tuple represents exactly the concrete child
configurations admissible for that e-node.

Read patiently, the requirement on \(\chi(n)\) says that \(n\)
appears in the relation once for every combination of members drawn
from its referenced classes, and for no others: if \(n\) admits
child alternatives slot-wise, it admits every mixture of them, and
the child positions never interact.  We call this the
\emph{Cartesian-fiber condition}.  Combinatorialists call a
Cartesian product of sets a \emph{rectangle}, after the
two-dimensional case, in which the tuples of \(C_1\times C_2\) form
a solid block in the matrix of the relation; in that vocabulary, the
condition says the inverse image of each e-node is one full
rectangle.  Both vocabularies earn their keep later, when e-nodes
are counted (Remark~\ref{rem:rectangles}).

If a concrete composition fiber cannot be expressed as one Cartesian
product of multe-classes, it is represented by multiple e-nodes, each
with its own child-reference tuple.

\begin{definition}[Multe-graph]
A \emph{multe-graph} is a tuple
\[
\mathcal G=(\mathcal S,\mathcal C,\chi)
\]
consisting of

\begin{enumerate}

\item
a concrete semantic schema
\(\mathcal S=(X,\mathsf N,\Sigma,\mathcal R)\)
(Definition~\ref{def:concrete-schema});

\item
a refinement hierarchy of multe-classes
\(\mathcal C(x)\) for each goal \(x\); and

\item
a child-reference tuple \(\chi(n)\) satisfying the condition above for
every non-atomic e-node \(n\).

\end{enumerate}
\end{definition}

A multe-graph contains three orthogonal relations:

\begin{itemize}

\item
\emph{membership}: multe-classes contain e-nodes;

\item
\emph{refinement}: finer multe-classes are contained in coarser ones;
and

\item
\emph{composition}: non-atomic e-nodes reference multe-classes of
subgoals.

\end{itemize}

The child reference \(C_i\) records precisely the distinctions among
the \(i\)-th subgoal realizations that the parent e-node observes.
Any member of \(C_i\) may be used, and no realization outside \(C_i\)
may be used.

\begin{example}[Join alternatives]
Let
\[
d_{\mathsf{join}}
=
(x_{R\Join S}\Rightarrow x_R,x_S).
\]

Suppose the input hierarchies contain
\[
[R]_{\mathsf{sorted}(a)}
\subseteq
[R]_{\mathsf{any}}
\]
and
\[
[S]_{\mathsf{sorted}(a)}
\subseteq
[S]_{\mathsf{any}}.
\]

A hash-join e-node may have child references
\[
\chi(n_{\mathsf{hash}})
=
\bigl(
[R]_{\mathsf{any}},
[S]_{\mathsf{any}}
\bigr),
\]
while a merge-join e-node may have
\[
\chi(n_{\mathsf{merge}})
=
\bigl(
[R]_{\mathsf{sorted}(a)},
[S]_{\mathsf{sorted}(a)}
\bigr).
\]

Neither e-node selects particular implementations of \(R\) or \(S\).
Each identifies exactly the classes of implementations admissible at
its child positions.
\end{example}

Figure~\ref{fig:selinger-multe-graph} displays the refinement and
composition structures together.  Nesting records refinement among
multe-classes, while arrows record the child references of the
non-atomic join e-nodes.

\begin{figure}[t]
\centering
\resizebox{\columnwidth}{!}{%
\begin{tikzpicture}[
  enode/.style={draw, rounded corners=2pt, fill=white,
    inner xsep=5pt, inner ysep=3pt, font=\small},
  mclass/.style={draw, rounded corners=5pt, inner sep=5pt},
  sortedbox/.style={draw, rounded corners=4pt, inner sep=4pt},
  classlabel/.style={fill=white, inner sep=1pt, font=\small},
  goal/.style={font=\small\bfseries},
  ref/.style={-{Latex[length=2.2mm]}, semithick},
  anyref/.style={ref, dashed},
  sortedref/.style={ref}
]

% ===== Parent goal x_{R Join S} =====
\node[enode] (merge) at (-1.4,0) {$\mathsf{MergeJoin}_a$};
\node[enode, right=9mm of merge] (hash) {$\mathsf{HashJoin}$};
\node[classlabel, anchor=south west] (RSsortedlab)
  at ([xshift=-3pt,yshift=2.2mm]merge.north west)
  {$[R\Join S]_{\mathsf{sorted}(a)}$};
\node[sortedbox, fit=(merge)(RSsortedlab)] (RSsorted) {};
\node[classlabel, anchor=south west] (RSanylab)
  at ([xshift=-3pt,yshift=2.2mm]RSsorted.north west)
  {$[R\Join S]_{\mathsf{any}}$};
\node[mclass, fit=(RSsorted)(hash)(RSanylab)] (RSany) {};
\node[goal, above=2.5mm of RSany] {$x_{R\Join S}$};

% ===== Child goal x_R =====
\node[enode] (rindex) at (-4.7,-4.0) {$\mathsf{IndexScan}_a$};
\node[enode, right=8mm of rindex] (rtable) {$\mathsf{TableScan}$};
\node[classlabel, anchor=south west] (Rsortedlab)
  at ([xshift=-3pt,yshift=2.2mm]rindex.north west)
  {$[R]_{\mathsf{sorted}(a)}$};
\node[sortedbox, fit=(rindex)(Rsortedlab)] (Rsorted) {};
\node[classlabel, anchor=south west] (Ranylab)
  at ([xshift=-3pt,yshift=2.2mm]Rsorted.north west)
  {$[R]_{\mathsf{any}}$};
\node[mclass, fit=(Rsorted)(rtable)(Ranylab)] (Rany) {};
\node[goal, below=2.5mm of Rany] {$x_R$};

% ===== Child goal x_S =====
\node[enode] (sindex) at (1.9,-4.0) {$\mathsf{IndexScan}_a$};
\node[enode, right=8mm of sindex] (stable) {$\mathsf{TableScan}$};
\node[classlabel, anchor=south west] (Ssortedlab)
  at ([xshift=-3pt,yshift=2.2mm]sindex.north west)
  {$[S]_{\mathsf{sorted}(a)}$};
\node[sortedbox, fit=(sindex)(Ssortedlab)] (Ssorted) {};
\node[classlabel, anchor=south west] (Sanylab)
  at ([xshift=-3pt,yshift=2.2mm]Ssorted.north west)
  {$[S]_{\mathsf{any}}$};
\node[mclass, fit=(Ssorted)(stable)(Sanylab)] (Sany) {};
\node[goal, below=2.5mm of Sany] {$x_S$};

% ===== Child-reference functions =====
\draw[sortedref] (merge.south) to[bend right=12] (Rsorted.north);
\draw[sortedref] (merge.south) to[bend left=22]  (Ssorted.north);
\draw[anyref]    (hash.south)  to[bend right=24] (Rany.north east);
\draw[anyref]    (hash.south)  to[bend left=10]  (Sany.north west);

\end{tikzpicture}%
}

\caption{
A multe-graph for a Selinger-style join.
Nested boxes are multe-classes and rounded interior boxes are
e-nodes.
The merge-join node references the sorted input classes, while
the hash-join node references the coarser input classes.
Solid arrows denote the sorted references; dashed arrows denote
the coarser references.
}
\Description{A multe-graph with three nested-box clusters: a parent
join goal above two child goals for the base relations R and S. Each
cluster is an outer ``any'' multe-class box containing a nested
``sorted(a)'' multe-class box. The parent's sorted box holds a
merge-join e-node and its any box also holds a hash-join e-node; each
child's sorted box holds an index-scan e-node and its any box also
holds a table-scan e-node. Solid arrows run from the merge-join node
to the two children's sorted boxes; dashed arrows run from the
hash-join node to the two children's any boxes.}
\label{fig:selinger-multe-graph}
\end{figure}

%%%%%%%%%%%%%%%%%%%%%%%%%%%%%%%%%%%%%%%%%%%%%%%%%%%%%%%%%%%%%%%
\section{Problem Instances}
\label{sec:instances}
%%%%%%%%%%%%%%%%%%%%%%%%%%%%%%%%%%%%%%%%%%%%%%%%%%%%%%%%%%%%%%%

A multe-graph is a schema for a family of related problems.
It records the available e-nodes, their composition, and the semantic
distinctions represented by the multe-class hierarchies.
A particular problem instance is obtained by registering
\emph{requirements} with nodes of that schema.

We begin with Boolean requirements.
More expressive requirement domains will be considered later.

%%%%%%%%%%%%%%%%%%%%%%%%%%%%%%%%%%%%%%%%%%%%%%%%%%%%%%%%%%%%%%%
\subsection{Requirement Predicates}
\label{sec:requirement-predicates}
%%%%%%%%%%%%%%%%%%%%%%%%%%%%%%%%%%%%%%%%%%%%%%%%%%%%%%%%%%%%%%%

Let \(C\in\mathcal C(x)\) be a multe-class for some goal \(x\).

\begin{definition}[Requirement]
A requirement on \(C\) is a predicate
\[
\rho:C\longrightarrow\{\mathsf{true},\mathsf{false}\}.
\]

The e-nodes in \(C\) that satisfy the requirement are
\[
C\!\restriction_{\rho}
=
\{n\in C\mid \rho(n)=\mathsf{true}\}.
\]
\end{definition}

In this section we restrict attention to \emph{local requirements}.
A local requirement can be evaluated from the e-node itself:
\(\rho(n)\) does not depend on which realizations are selected
beneath \(n\).

The multe-class \(C\) and the requirement \(\rho\) play different
roles.

The class \(C\) is part of the semantic schema.
It records a set of e-nodes that are interchangeable under one
represented semantic equivalence.

The requirement \(\rho\) is part of a problem instance.
It restricts the members of \(C\) that remain admissible for that
instance.

A requirement need not coincide with a represented multe-class.
A requirement on \(C\) may select any subset of \(C\).
When a desired subset is already represented by a nested
multe-class, however, that class induces a requirement in the obvious
way.

\begin{definition}[Class-induced requirement]
Let \(D,C\in\mathcal C(x)\) with \(D\subseteq C\).
The requirement on \(C\) induced by \(D\) is the characteristic
predicate
\(\rho_D(n) \Leftrightarrow n\in D\).
Its satisfying set is exactly \(C\!\restriction_{\rho_D}=D\).
\end{definition}

Thus selecting a multe-class and imposing a requirement are not
different mechanisms.
Selecting a represented subclass is one particularly simple
requirement.

Requirements induced by represented multe-classes are local.  For
example, whether a root e-node belongs to
\([R\Join S]_{\mathsf{sorted}(a)}\) is intrinsic to that e-node's
provided semantics.

Several requirements may be registered at the same class.
Their conjunction, \((\rho_1\wedge\rho_2)(n) =
\rho_1(n)\wedge\rho_2(n)\), is again a requirement, so we may treat
any finite collection of requirements at one site as a single
predicate.

Requirements may be imposed on multiple multe-classes in a graph.
A requirement on an internal class may, for example, rule out sorting
implementations, require a particular execution device, or impose a
resource restriction.
In the next subsection we begin with the simplest case: a single
requirement registered at the root.

%%%%%%%%%%%%%%%%%%%%%%%%%%%%%%%%%%%%%%%%%%%%%%%%%%%%%%%%%%%%%%%
\subsection{Root Requirements}
\label{sec:root-requirements}
%%%%%%%%%%%%%%%%%%%%%%%%%%%%%%%%%%%%%%%%%%%%%%%%%%%%%%%%%%%%%%%

A problem instance designates one goal as the result to be realized.
Let \(x_\star\in X\) be that root goal.  Its refinement hierarchy
has the universal root class
\(C^\top_\star = \mathsf N(x_\star)\).

\begin{definition}[Root requirement]
A \emph{root requirement} is a predicate
\[
\rho_\star:
C^\top_\star
\longrightarrow
\{\mathsf{true},\mathsf{false}\}
\]
registered at the universal multe-class of the root goal.
\end{definition}

The satisfying set \(C^\top_\star\!\restriction_{\rho_\star}\)
contains the root e-nodes that provide an acceptable result for the
instance.

\begin{definition}[Root-only problem instance]
A \emph{root-only problem instance} is a tuple
\[
I=(\mathcal G,x_\star,\rho_\star),
\]
where

\begin{enumerate}

\item
\(\mathcal G\) is a multe-graph;

\item
\(x_\star\) is the root goal; and

\item
\(\rho_\star\) is a requirement registered at
\(C^\top_\star=\mathsf N(x_\star)\).

\end{enumerate}
\end{definition}

The predicate that is identically true expresses no additional
restriction:
\[
\rho_{\mathsf{any}}(n)=\mathsf{true}
\qquad
\text{for every }n\in C^\top_\star.
\]

A more selective root requirement may be induced by a nested
multe-class.

\begin{example}[The running join queries]
Let \(x_\star=x_{R\Join S}\),
and suppose the root hierarchy contains
\[
[R\Join S]_{\mathsf{sorted}(a)}
\subseteq
[R\Join S]_{\mathsf{any}},
\]
where
\[
[R\Join S]_{\mathsf{any}}
=
\mathsf N(x_{R\Join S}).
\]

The query
\[
\texttt{SELECT * FROM R JOIN S}
\]
imposes the root requirement
\(\rho_{\mathsf{any}}(n)=\mathsf{true}\).

The query
\[
\texttt{SELECT * FROM R JOIN S ORDER BY R.a}
\]
imposes the class-induced root requirement
\[
\rho_{\mathsf{sorted}(a)}(n)
\quad\Longleftrightarrow\quad
n\in[R\Join S]_{\mathsf{sorted}(a)}.
\]

The two queries therefore share one multe-graph.
They differ only in the predicate registered at its root.
\end{example}

This formulation also permits root requirements that do not correspond
to a pre-existing multe-class.
For example, a predicate in a distributed database might only accept root e-nodes with a ``location'' label set to the client location.
Whether such a requirement should itself be added to the schema as a
new represented semantic class is a separate design question.

A root requirement restricts the admissible root e-nodes, but its
consequences are not confined to the root.
The surviving root e-nodes reference child multe-classes; requirements
on those classes may in turn affect their parents.
The next subsection defines how registered requirements propagate
through the multe-graph and produce a problem-specific instantiated
graph.

Before we get there, note that not every useful requirement is local.
An \emph{aggregate requirement} depends on values computed from the
realizations selected for an e-node's children.  Total cost, latency,
memory consumption, and energy use are typical examples.

Aggregate requirements require additional compositional structure:
each non-atomic e-node must specify how child values are combined into
a parent value.  Propagating such a requirement then requires
distributing it through that aggregation function to obtain
requirements or bounds on the children.  We postpone this extension
until we introduce cost-based optimization, where it will also support
branch-and-bound pruning.

%%%%%%%%%%%%%%%%%%%%%%%%%%%%%%%%%%%%%%%%%%%%%%%%%%%%%%%%%%%%%%%
\subsection{Propagation and Instantiation}
\label{sec:propagation}
%%%%%%%%%%%%%%%%%%%%%%%%%%%%%%%%%%%%%%%%%%%%%%%%%%%%%%%%%%%%%%%

A root requirement identifies acceptable root e-nodes.
Not every such e-node necessarily participates in a complete
realization: a non-atomic e-node is usable only if each of its
referenced child multe-classes contains a usable e-node.

Propagation therefore has two directions.

\begin{enumerate}

\item
\emph{Realizability} is computed upward from atomic e-nodes.

\item
\emph{Liveness} is propagated downward from the root e-nodes that
satisfy the requirement.

\end{enumerate}

We continue to assume that the multe-graph is finite and acyclic, so
both notions can be defined inductively.

%%%%%%%%%%%%%%%%%%%%%%%%%%%%%%%%%%%%%%%%%%%%%%%%%%%%%%%%%%%%%%%
\subsubsection{Realizability}
%%%%%%%%%%%%%%%%%%%%%%%%%%%%%%%%%%%%%%%%%%%%%%%%%%%%%%%%%%%%%%%

\begin{definition}[Realizable e-node]
An atomic e-node is \emph{realizable}.

A non-atomic e-node \(n\), with child references
\(\chi(n)=(C_1,\ldots,C_k)\),
is realizable when every referenced child class contains at least one
realizable e-node.
\end{definition}

For a goal \(x\), let
\(\mathsf{Real}(x) =
\{n\in\mathsf N(x)\mid n\text{ is realizable}\}\),
and for a multe-class \(C\in\mathcal C(x)\), let
\(\mathsf{Real}(C) = C\cap\mathsf{Real}(x)\).
Thus \(n\) is realizable exactly when
\(\mathsf{Real}(C_i)\neq\varnothing\) for every child reference
\(C_i\).

Realizability does not depend on the root requirement.  It removes
dead ends already present in the schema.

%%%%%%%%%%%%%%%%%%%%%%%%%%%%%%%%%%%%%%%%%%%%%%%%%%%%%%%%%%%%%%%
\subsubsection{Root Candidates and Liveness}
%%%%%%%%%%%%%%%%%%%%%%%%%%%%%%%%%%%%%%%%%%%%%%%%%%%%%%%%%%%%%%%

Let \(I=(\mathcal G,x_\star,\rho_\star)\)
be a root-only problem instance.
The realizable root e-nodes satisfying the root requirement are
\[
\mathsf{Root}(I)
=
\left\{
n\in\mathsf{Real}(x_\star)
\;\middle|\;
\rho_\star(n)=\mathsf{true}
\right\}.
\]

If \(\mathsf{Root}(I)=\varnothing\),
then the problem instance is unsatisfiable.
Otherwise, these e-nodes initiate downward propagation.

\begin{definition}[Live e-node]
The set of \emph{live e-nodes} for \(I\), written
\(\mathsf{Live}(I)\),
is the least set satisfying the following rules.

\begin{enumerate}

\item
Every acceptable realizable root e-node is live:
\(\mathsf{Root}(I)\subseteq\mathsf{Live}(I)\).

\item
If \(n\in\mathsf{Live}(I)\) and
\(\chi(n)=(C_1,\ldots,C_k)\),
then every realizable e-node in every referenced child class is live:
\[
\mathsf{Real}(C_i)\subseteq\mathsf{Live}(I)
\qquad
\text{for each }i.
\]

\end{enumerate}
\end{definition}

This closure retains every realizable alternative that may contribute
to a realization satisfying the root requirement.  It does not choose
one alternative from any class.

For each goal \(x\), define its live e-node set by
\(\mathsf N_I(x) = \mathsf N(x)\cap\mathsf{Live}(I)\).

A multe-class is \emph{referenced in \(I\)} if it occurs in
\(\chi(n)\) for some live non-atomic e-node \(n\).

%%%%%%%%%%%%%%%%%%%%%%%%%%%%%%%%%%%%%%%%%%%%%%%%%%%%%%%%%%%%%%%
\subsubsection{The Instantiated Multe-Graph}
%%%%%%%%%%%%%%%%%%%%%%%%%%%%%%%%%%%%%%%%%%%%%%%%%%%%%%%%%%%%%%%

We obtain the problem-specific graph by restricting both e-nodes and
multe-classes to the live portion of the schema.

For every referenced multe-class \(C\in\mathcal C(x)\),
define its restriction to \(I\) by
\(C\!\restriction_I = C\cap\mathsf N_I(x)\).
The active universe \(\mathsf N_I(x)\)
serves as the root multe-class for each live goal \(x\).

\begin{definition}[Instantiated multe-graph]
The \emph{instantiated multe-graph} for \(I\), written
\(\mathcal G\!\restriction_I\),
is obtained as follows.

\begin{enumerate}

\item
For each goal \(x\), retain the live e-nodes \(\mathsf N_I(x)\).

\item
Retain the nonempty restrictions \(C\!\restriction_I\)
of all multe-classes referenced by live e-nodes, together with the
active root class \(\mathsf N_I(x)\).

\item
Identify restricted classes that have become equal, and order the
remaining classes by inclusion.

\item
For every live non-atomic e-node \(n\) with
\(\chi(n)=(C_1,\ldots,C_k)\),
replace its child references by
\[
\chi_I(n)
=
\bigl(
C_1\!\restriction_I,\ldots,
C_k\!\restriction_I
\bigr).
\]

\end{enumerate}
\end{definition}

Because the original multe-classes form a tree, any two classes are
\emph{laminar}: either nested or disjoint.  Intersecting them with the same live
e-node set preserves the \emph{laminarity} property: if
\(C\subseteq D\), then
\(C\cap\mathsf N_I(x) \subseteq D\cap\mathsf N_I(x)\),
and if \(C\cap D=\varnothing\), their restrictions remain disjoint.

Identifying equal restrictions removes boundaries that no longer
distinguish any live e-nodes, so once empty restrictions are discarded
and equal ones identified, the remaining classes again form a
tree-shaped refinement hierarchy.

Classes that are neither selected at the root nor referenced by a live
e-node do not appear in the instantiated graph.  Their distinctions
cannot be observed by any realization satisfying the root requirement.

This is the first sense in which a multe-graph is distilled.
The original graph records every distinction that may matter to some
problem instance; the instantiated graph retains only distinctions
reachable through realizable alternatives from the current root
requirement.

Because a live e-node retains every realizable member of each
referenced child class, its instantiated child space remains a
Cartesian product:
\[
\mathsf{Real}(C_1)
\times\cdots\times
\mathsf{Real}(C_k).
\]
Thus instantiation preserves the well-formedness condition built into
the definition of a multe-graph.

\begin{example}[The running join instances]
Suppose the root hierarchy contains
\[
[R\Join S]_{\mathsf{sorted}(a)}
\subseteq
[R\Join S]_{\mathsf{any}}.
\]

For the query
\[
\texttt{SELECT * FROM R JOIN S},
\]
the root requirement accepts every root e-node.
If both hash join and merge join are realizable, both become live.
Their child references make the coarse and sorted input classes live,
respectively.

For the query
\[
\texttt{SELECT * FROM R JOIN S ORDER BY R.a},
\]
the root requirement accepts only e-nodes in
\([R\Join S]_{\mathsf{sorted}(a)}\).
If merge join is the only realizable e-node in that class, then its
references make
\[
[R]_{\mathsf{sorted}(a)}
\qquad\text{and}\qquad
[S]_{\mathsf{sorted}(a)}
\]
live.

In either case, all realizable e-nodes in each referenced input class
remain available.  Instantiation prunes unreachable alternatives but
does not yet select representatives.
\end{example}

Instantiation determines which alternatives and semantic distinctions
survive the root requirement.  Choosing among the surviving
alternatives is the separate task of canonicalization.

%%%%%%%%%%%%%%%%%%%%%%%%%%%%%%%%%%%%%%%%%%%%%%%%%%%%%%%%%%%%%%%
\section{Canonicalization}
\label{sec:canonicalization}
%%%%%%%%%%%%%%%%%%%%%%%%%%%%%%%%%%%%%%%%%%%%%%%%%%%%%%%%%%%%%%%

Instantiation removes alternatives that cannot contribute to a
realization satisfying the problem requirements.  It may nevertheless
leave many e-nodes in each surviving multe-class.

Canonicalization performs a separate task: it chooses how each
surviving class will be represented.

In the setting considered here, canonicalization represents every
surviving multe-class by exactly one of its e-nodes.

%%%%%%%%%%%%%%%%%%%%%%%%%%%%%%%%%%%%%%%%%%%%%%%%%%%%%%%%%%%%%%%
\subsection{Canonicalization Policies}
\label{sec:canonicalization-policies}
%%%%%%%%%%%%%%%%%%%%%%%%%%%%%%%%%%%%%%%%%%%%%%%%%%%%%%%%%%%%%%%

Let \(\mathcal G\!\restriction_I\)
be the instantiated multe-graph for a satisfiable problem instance
\(I\).  For each goal \(x\), let \(\mathcal C_I(x)\)
denote the nonempty multe-classes retained in the instantiated
graph, and let
\(\mathcal C_I = \bigcup_{x\in X}\mathcal C_I(x)\).

\begin{definition}[Canonicalization]
A \emph{canonicalization} of
\(\mathcal G\!\restriction_I\) is an assignment \(\kappa\)
that selects one e-node from every retained multe-class:
\[
\kappa(C)\in C
\qquad
\text{for every }C\in\mathcal C_I.
\]

The selected e-node \(\kappa(C)\) is the
\emph{canonical representative} of \(C\).
\end{definition}

Because every retained class is nonempty, a canonicalization always
exists in the finite setting considered here.  The definition does not
yet specify which member should be selected.

\begin{definition}[Canonicalization policy]
A \emph{canonicalization policy} is a rule that produces a
canonicalization for each instantiated multe-graph in its domain.
\end{definition}

The distinction between requirements and canonicalization is now
precise.

A requirement restricts the admissible members of a class:
\(C \mapsto C\!\restriction_{\rho}\).
Canonicalization acts only after those restrictions have been
propagated and selects one of the surviving members:
\(C\!\restriction_{\rho} \mapsto
\kappa(C\!\restriction_{\rho})\).

A requirement may leave several admissible alternatives.
Canonicalization deliberately removes that remaining multiplicity.

Canonicalization does not alter the semantic meaning of a
multe-class.  Since \(\kappa(C)\in C\),
the selected e-node provides the semantics represented by \(C\).
The policy determines which such provider is retained.

Selections for nested multe-classes are independent.  If
\(C\subseteq D\), there is no requirement that
\(\kappa(C)=\kappa(D)\).
The two classes represent different semantic contracts and may have
different preferred representatives.

\begin{example}[Canonical representatives for the running join]
Suppose an instantiated graph contains
\[
[R\Join S]_{\mathsf{sorted}(a)}
\subseteq
[R\Join S]_{\mathsf{any}}.
\]

A canonicalization may select
\(\kappa\bigl([R\Join S]_{\mathsf{any}}\bigr) = n_{\mathsf{hash}}\)
and
\(\kappa\bigl([R\Join S]_{\mathsf{sorted}(a)}\bigr) =
n_{\mathsf{merge}}\).
There is no conflict between these choices.
The hash join represents the coarser class of arbitrary join plans,
while the merge join represents the finer class of plans that provide
the required order.
\end{example}

A selected non-atomic e-node continues to reference child
multe-classes.  Canonicalization supplies one representative for each
of those classes as well.  Thus, beginning from a root class and
repeatedly following the selected e-nodes and their child references
determines a complete realization, and because the current
multe-graphs are acyclic, this recursive process terminates at atomic
e-nodes.

Canonicalization therefore reduces the instantiated multe-graph to one
compositional choice per surviving semantic class.  It does not,
however, explain why those choices are preferable.

An arbitrary policy may select arbitrary members.  An optimization
policy must instead compare complete realizations according to an
objective such as cost.  Since the value of a non-atomic realization
generally depends on the values of its children, defining such a policy
requires a compositional account of aggregate attributes.  We introduce
that structure next.

%%%%%%%%%%%%%%%%%%%%%%%%%%%%%%%%%%%%%%%%%%%%%%%%%%%%%%%%%%%%%%%
\section{Extraction}
\label{sec:extraction}
%%%%%%%%%%%%%%%%%%%%%%%%%%%%%%%%%%%%%%%%%%%%%%%%%%%%%%%%%%%%%%%

Instantiation and canonicalization together determine a complete
realization of the root goal.

Let \(I=(\mathcal G,x_\star,\rho_\star)\)
be a satisfiable problem instance, let
\(\mathcal G\!\restriction_I\)
be its instantiated multe-graph, and let \(\kappa\)
be a canonicalization of that graph.
The active root class is
\(C^\star_I=\mathsf N_I(x_\star)\).
By construction, every e-node in \(C^\star_I\) is realizable and
satisfies the root requirement.

%%%%%%%%%%%%%%%%%%%%%%%%%%%%%%%%%%%%%%%%%%%%%%%%%%%%%%%%%%%%%%%
\subsection{Extracted Realizations}
\label{sec:extracted-realizations}
%%%%%%%%%%%%%%%%%%%%%%%%%%%%%%%%%%%%%%%%%%%%%%%%%%%%%%%%%%%%%%%

Canonicalization selects an e-node from each retained multe-class.
The selected e-node may itself reference child multe-classes, whose
selected e-nodes must in turn be followed.

\begin{definition}[Extracted realization]
Let \(C\) be a retained multe-class, and let \(n=\kappa(C)\).
The \emph{realization extracted from \(C\)}, written
\(\mathsf{Extract}_{\kappa}(C)\),
is defined recursively as follows.
If \(n\) is atomic, then
\(\mathsf{Extract}_{\kappa}(C) = \langle n\rangle\).
If \(n\) is non-atomic and has instantiated child references
\(\chi_I(n)=(C_1,\ldots,C_k)\),
then
\[
\mathsf{Extract}_{\kappa}(C)
=
\left\langle
n,
\mathsf{Extract}_{\kappa}(C_1),
\ldots,
\mathsf{Extract}_{\kappa}(C_k)
\right\rangle.
\]
\end{definition}

Because the instantiated multe-graph is acyclic, this recursion
terminates at atomic e-nodes.

The extracted realization may be viewed as a tree.  If repeated
references to the same multe-class share their extracted
realization, it may instead be represented as a directed acyclic
graph.  This distinction affects representation but not semantics.

\begin{definition}[Extraction]
Given a problem instance \(I\) and a canonicalization \(\kappa\) of
its instantiated multe-graph, the \emph{extraction result} is
\[
\mathsf{Extract}(I,\kappa)
=
\mathsf{Extract}_{\kappa}(C^\star_I).
\]
\end{definition}

Thus extraction has three conceptual steps:

\begin{enumerate}

\item
propagate the requirements to obtain
\(\mathcal G\!\restriction_I\);

\item
canonicalize each retained multe-class; and

\item
follow the canonical choices recursively from the active root class.

\end{enumerate}

No evaluation order or search algorithm is part of this definition.

%%%%%%%%%%%%%%%%%%%%%%%%%%%%%%%%%%%%%%%%%%%%%%%%%%%%%%%%%%%%%%%
\subsection{Correctness}
\label{sec:extraction-correctness}
%%%%%%%%%%%%%%%%%%%%%%%%%%%%%%%%%%%%%%%%%%%%%%%%%%%%%%%%%%%%%%%

\begin{proposition}[Validity of extraction]
For every retained multe-class \(C\),
\(\mathsf{Extract}_{\kappa}(C)\)
is a concrete realization of the goal associated with \(C\), rooted at
an e-node belonging to \(C\).
In particular, \(\mathsf{Extract}(I,\kappa)\)
is a concrete realization of \(x_\star\) satisfying the root
requirement \(\rho_\star\).
\end{proposition}

\begin{proof}
The claim follows by induction over the acyclic composition graph.
If \(\kappa(C)=n\)
is atomic, then \(n\in C\) is itself a realization of the goal
associated with \(C\).
Otherwise, let \(\chi_I(n)=(C_1,\ldots,C_k)\).
By induction, each \(\mathsf{Extract}_{\kappa}(C_i)\)
is rooted at an e-node belonging to \(C_i\).
The child-reference condition for the multe-graph therefore ensures
that these child realizations form an admissible concrete composition
with \(n\).
Since \(n=\kappa(C)\in C\), the resulting realization is represented
by \(C\).

At the root, \(\kappa(C^\star_I)\in C^\star_I\),
and every member of \(C^\star_I\) satisfies \(\rho_\star\).
\end{proof}

The proposition establishes only validity: an arbitrary
canonicalization produces an arbitrary valid extraction.  To establish
optimality, we must define an aggregate objective and show that local
canonical choices compose into a globally optimal realization.  That is
the role of compositional cost and Bellman compatibility.

%%%%%%%%%%%%%%%%%%%%%%%%%%%%%%%%%%%%%%%%%%%%%%%%%%%%%%%%%%%%%%%
\section{Compositional Attributes and Optimization}
\label{sec:optimization}
%%%%%%%%%%%%%%%%%%%%%%%%%%%%%%%%%%%%%%%%%%%%%%%%%%%%%%%%%%%%%%%

The preceding sections define valid extraction independently of any
objective.  We now add attributes whose values depend compositionally
on complete realizations.  Cost is the principal example.

%%%%%%%%%%%%%%%%%%%%%%%%%%%%%%%%%%%%%%%%%%%%%%%%%%%%%%%%%%%%%%%
\subsection{Compositional Cost}
\label{sec:compositional-cost}
%%%%%%%%%%%%%%%%%%%%%%%%%%%%%%%%%%%%%%%%%%%%%%%%%%%%%%%%%%%%%%%

Canonicalization selects an e-node for each surviving multe-class.
To choose among those e-nodes by cost, we must assign a cost to each
complete realization.

We begin with the simplest familiar model: every e-node has a
nonnegative local cost, and the cost of a realization is the sum of the
local costs of its e-nodes.

\begin{definition}[Local cost]
A \emph{local cost function} is a function
\(c:\mathsf N\to\mathbb R_{\geq 0}\)
that assigns a nonnegative cost \(c(n)\) to every e-node \(n\).
\end{definition}

The local cost of a non-atomic e-node does not include the costs of its
children.  Those costs are incorporated recursively.

\begin{definition}[Cost of a realization]
Let
\[
t=
\left\langle
n,t_1,\ldots,t_k
\right\rangle
\]
be a complete realization rooted at an e-node \(n\).
Its cost is
\[
\operatorname{cost}(t)
=
c(n)
+
\sum_{i=1}^{k}\operatorname{cost}(t_i).
\]

For an atomic realization \(t=\langle n\rangle\),
\(\operatorname{cost}(t)=c(n)\).
\end{definition}

Cost is therefore an aggregate property of the complete realization.
It cannot generally be determined from the root e-node alone.

A canonicalization determines one complete realization below every
retained multe-class and therefore induces a cost for that class.

Let \(\kappa(C)=n\),
and suppose that the instantiated child references of \(n\) are
\(\chi_I(n)=(C_1,\ldots,C_k)\).

\begin{definition}[Cost induced by a canonicalization]
The cost induced by \(\kappa\) at \(C\), written
\(\operatorname{cost}_{\kappa}(C)\),
is defined recursively by
\[
\operatorname{cost}_{\kappa}(C)
=
c(n)
+
\sum_{i=1}^{k}
\operatorname{cost}_{\kappa}(C_i),
\qquad
n=\kappa(C).
\]

If \(\kappa(C)\) is atomic, then
\[
\operatorname{cost}_{\kappa}(C)
=
c\bigl(\kappa(C)\bigr).
\]
\end{definition}

Because the instantiated multe-graph is acyclic, this recursion is
well-defined.  Equivalently,
\[
\operatorname{cost}_{\kappa}(C)
=
\operatorname{cost}
\bigl(
\mathsf{Extract}_{\kappa}(C)
\bigr).
\]

\begin{example}[Cost of the running join]
Suppose
\[
\kappa\bigl([R\Join S]_{\mathsf{any}}\bigr)
=
n_{\mathsf{hash}},
\]
and the hash-join e-node references
\[
[R]_{\mathsf{any}}
\qquad\text{and}\qquad
[S]_{\mathsf{any}}.
\]

Then
\[
\begin{split}
\operatorname{cost}_{\kappa}
\bigl([R\Join S]_{\mathsf{any}}\bigr)
={}&
c(n_{\mathsf{hash}})
\\
&+
\operatorname{cost}_{\kappa}
\bigl([R]_{\mathsf{any}}\bigr)
\\
&+
\operatorname{cost}_{\kappa}
\bigl([S]_{\mathsf{any}}\bigr).
\end{split}
\]

Thus the cost of choosing the hash-join e-node depends on the
canonical realizations chosen for its referenced child classes.
\end{example}

This additive model supports two distinct uses of cost.

First, cost may appear in an aggregate requirement, such as
\(\operatorname{cost}(t)\leq B\).
Unlike the local requirements considered earlier, such a predicate
depends on the recursively selected children and must be propagated
through the cost equation.

Second, cost may define a canonicalization policy that selects a
minimum-cost realization from each surviving multe-class.

We consider cost requirements and their propagation next.  More general
aggregation functions, including non-additive and operator-specific
ones, are deferred to Section~\ref{sec:generalizations}.

%%%%%%%%%%%%%%%%%%%%%%%%%%%%%%%%%%%%%%%%%%%%%%%%%%%%%%%%%%%%%%%
\subsection{Cost Requirements and Propagation}
\label{sec:cost-requirements}
%%%%%%%%%%%%%%%%%%%%%%%%%%%%%%%%%%%%%%%%%%%%%%%%%%%%%%%%%%%%%%%

The requirements considered in
Section~\ref{sec:requirement-predicates} are \emph{local}: whether an
e-node satisfies them can be determined from the e-node itself.

A cost requirement is different.  It applies to a complete realization,
and its truth depends on the realizations selected for the children.
It must therefore be propagated through the compositional cost
equation.  Let \(\mathcal T_I(C)\) denote the set of complete
realizations represented by a retained multe-class \(C\) in the
instantiated graph \(\mathcal G\!\restriction_I\).

\begin{definition}[Cost requirement]
For a budget \(B\in\mathbb R_{\geq 0}\), the \emph{cost requirement}
with budget \(B\) is the predicate
\(\rho_B(t)\Leftrightarrow\operatorname{cost}(t)\leq B\) on complete
realizations.  A \emph{budgeted demand} is a pair \((C,B)\)
requiring a realization represented by \(C\) whose cost is at most
\(B\).
\end{definition}

Suppose an e-node \(n\in C\) has child references
\(\chi_I(n)=(C_1,\ldots,C_k)\).  A realization rooted at \(n\)
satisfies budget \(B\) exactly when its children satisfy
\[
c(n)+\sum_{i=1}^{k}\operatorname{cost}(t_i)\leq B,
\]
so the parent budget leaves \(B-c(n)\) to be divided among the
children.  In general, there is no unique such division.

\begin{definition}[Budget allocation]
For an e-node \(n\) of arity \(k\), the set of budget allocations
admitted by a parent budget \(B\) is
\[
\mathsf{Alloc}_n(B)
=
\left\{
(B_1,\ldots,B_k)\in\mathbb R_{\geq 0}^{k}
\;\middle|\;
c(n)+\sum_{i=1}^{k}B_i\leq B
\right\}.
\]
\end{definition}

An allocation \((B_1,\ldots,B_k)\in\mathsf{Alloc}_n(B)\) propagates
the parent demand \((C,B)\) to the child demands
\((C_1,B_1),\ldots,(C_k,B_k)\).  This propagation is relational: one
parent demand may admit many different allocations among the
children.

\begin{example}[Allocating a join budget]
Suppose a hash-join e-node has local cost
\(c(n_{\mathsf{hash}})=20\) and the root budget is \(B=100\).  Its
child budgets must satisfy \(B_R+B_S\leq 80\); possible allocations
include \((30,50)\), \((40,40)\), and \((60,20)\).  The parent
requirement does not determine one independent requirement for each
child.  It determines a relation among their permitted budgets.
\end{example}

We can now define the exact propagation semantics.

\begin{definition}[Satisfiable budgeted demand]
A budgeted demand \((C,B)\) is \emph{satisfiable} if one of the
following holds.

\begin{enumerate}

\item
There is an atomic e-node \(n\in C\) with \(c(n)\leq B\).

\item
There is a non-atomic e-node \(n\in C\) with child references
\(\chi_I(n)=(C_1,\ldots,C_k)\) and an allocation
\((B_1,\ldots,B_k)\in\mathsf{Alloc}_n(B)\) such that every child
demand \((C_i,B_i)\) is satisfiable.

\end{enumerate}
\end{definition}

Because the instantiated multe-graph is acyclic, this recursive
definition is well-founded.

\begin{proposition}[Exactness of budget propagation]
A budgeted demand \((C,B)\) is satisfiable if and only if there
exists a realization \(t\in\mathcal T_I(C)\) with
\(\operatorname{cost}(t)\leq B\).
\end{proposition}

\begin{proof}
The proof is by induction over the acyclic composition graph.  For
an atomic realization \(t=\langle n\rangle\), the claim follows from
\(\operatorname{cost}(t)=c(n)\).  For a non-atomic realization
\(t=\langle n,t_1,\ldots,t_k\rangle\), we have
\(\operatorname{cost}(t)=c(n)+\sum_{i}\operatorname{cost}(t_i)\).
If \(\operatorname{cost}(t)\leq B\), choose
\(B_i=\operatorname{cost}(t_i)\); then
\((B_1,\ldots,B_k)\in\mathsf{Alloc}_n(B)\), and each child demand is
satisfiable by induction.  Conversely, given a valid allocation and
satisfying child realizations, their aggregate cost is at most
\(c(n)+\sum_{i}B_i\leq B\).
\end{proof}

Local requirement propagation restricts the original multe-graph by
removing e-nodes and unreachable classes.  Cost propagation has a
different shape: it lifts each demanded class \(C\) to a family of
budget-indexed demands \((C,B)\).  Even for one parent e-node, exact
propagation may generate many alternative child-budget allocations;
the definition above supplies the semantics of those alternatives
without prescribing an algorithm for enumerating them.
Minimum-cost optimization provides a more direct characterization of
which budgets are satisfiable, and lower bounds will later allow
branch-and-bound algorithms to discard budget allocations and
e-nodes that cannot improve upon an incumbent solution.

%%%%%%%%%%%%%%%%%%%%%%%%%%%%%%%%%%%%%%%%%%%%%%%%%%%%%%%%%%%%%%%
\subsection{Minimum-Cost Canonicalization}
\label{sec:scalar-optimization}
%%%%%%%%%%%%%%%%%%%%%%%%%%%%%%%%%%%%%%%%%%%%%%%%%%%%%%%%%%%%%%%

The preceding section defined the cost of every complete realization.
We now use those costs to define a canonicalization policy.

For a retained multe-class \(C\), let \(\mathcal T_I(C)\)
denote the realizations represented by \(C\).

\begin{definition}[Optimal cost]
The \emph{optimal cost} of \(C\) is
\[
\operatorname{opt}(C)
=
\min_{t\in\mathcal T_I(C)}
\operatorname{cost}(t).
\]
\end{definition}

Because the instantiated multe-graph is finite and acyclic,
every retained multe-class represents finitely many realizations.
The minimum therefore exists.

A minimum-cost canonicalization chooses one realization attaining this
minimum.

\begin{definition}[Minimum-cost canonicalization]
A canonicalization \(\kappa\)
is \emph{minimum-cost} when, for every retained multe-class \(C\),
\[
\operatorname{cost}
\bigl(
\mathsf{Extract}_\kappa(C)
\bigr)
=
\operatorname{opt}(C).
\]
\end{definition}

Notice that this definition places no restriction on how \(\kappa\) is
computed; it merely specifies the desired outcome.  Nor does it assume
that optimal representatives may be chosen independently for different
multe-classes---indeed, without additional assumptions, choosing the
cheapest child representative need not produce the cheapest parent
realization.

The next subsection identifies the compatibility condition under which
local canonicalization is globally correct: Bellman's Principle of
Optimality.

%%%%%%%%%%%%%%%%%%%%%%%%%%%%%%%%%%%%%%%%%%%%%%%%%%%%%%%%%%%%%%%
\subsection{Bellman Compatibility}
\label{sec:bellman-compatibility}
%%%%%%%%%%%%%%%%%%%%%%%%%%%%%%%%%%%%%%%%%%%%%%%%%%%%%%%%%%%%%%%

The previous section defined minimum-cost canonicalization but did not
justify computing optimum realizations one multe-class at a time.  In a
general compositional optimization problem, such local reasoning is
unsound: the globally optimal realization of a parent need not be
assembled from locally optimal realizations of its children.  Dynamic
programming succeeds only when the semantic representation and the
optimization objective are compatible.

\begin{definition}[Bellman compatibility]
A compositional optimization problem is
\emph{Bellman-compatible} if the following conditions hold.

\begin{enumerate}

\item
\textbf{Semantic substitutability.}
Each child reference denotes a multe-class whose members are
semantically interchangeable in that position.

\item
\textbf{Independent composition.}
If an e-node references child classes \((C_1,\ldots,C_k)\), then
every tuple
\((t_1,\ldots,t_k)\in C_1\times\cdots\times C_k\)
is an admissible child configuration.  This is the Cartesian-fiber
condition of Section~\ref{sec:multe-graph-definition}, lifted from
e-nodes to the realizations they represent.

\item
\textbf{Monotone objective.}
Replacing a child realization by another realization of the same
multe-class whose cost is no greater cannot increase the cost of the
parent realization.

\end{enumerate}

\end{definition}

The first condition determines \emph{what} information is retained
about each child realization; the second ensures that children interact
only through that retained information; and the third ensures that
improving one child cannot harm its parent.  Together these conditions
imply the fundamental substitution property.

\begin{lemma}[Substitution]
Let an e-node reference child classes \((C_1,\ldots,C_k)\),
and suppose \(t_i,t_i'\in C_i\) with
\(\operatorname{cost}(t_i') \le \operatorname{cost}(t_i)\).
Replacing \(t_i\) by \(t_i'\) preserves admissibility and cannot
increase the parent's cost.
\end{lemma}

\begin{proof}
The replacement remains admissible because every realization of
\(C_i\) is interchangeable in that position and the children are
otherwise independent.  The cost claim follows from monotonicity.
\end{proof}

The classical Bellman principle is an immediate consequence.

\begin{theorem}[Bellman's Principle of Optimality]

In a Bellman-compatible optimization problem, an optimal realization of
every retained multe-class can be constructed from optimal realizations
of its referenced child classes.

Consequently, one optimal representative per retained multe-class
suffices for global optimization.

\end{theorem}

\paragraph{Discussion.}

Bellman's Principle is traditionally presented as a property of an
optimization objective: an optimal solution is assembled from optimal
solutions to subproblems.  The present development suggests a different
interpretation, in which the principle is not primitive but follows
from the interaction of a semantic representation and an objective.

The semantic representation supplies two ingredients: the refinement
hierarchy determines which distinctions among realizations are
preserved, and the child references determine the interfaces through
which subproblems interact.  The objective contributes only the
requirement that improving a child through such an interface cannot
worsen its parent.

Viewed this way, dynamic programming is not fundamentally about
recursion or memoization; it is optimization over semantic interfaces.
Different dynamic programming algorithms differ primarily in how they
choose those interfaces, and Bellman's Principle is the theorem that
those interfaces are sufficient for independent optimization.

%%%%%%%%%%%%%%%%%%%%%%%%%%%%%%%%%%%%%%%%%%%%%%%%%%%%%%%%%%%%%%%
\subsection{Bounds and Branch-and-Bound}
\label{sec:bounds-pruning}
%%%%%%%%%%%%%%%%%%%%%%%%%%%%%%%%%%%%%%%%%%%%%%%%%%%%%%%%%%%%%%%

Minimum-cost canonicalization is defined in terms of the exact
optimum cost \(\operatorname{opt}(C)\) of each retained multe-class.

Computing these exact values may itself require substantial search.
Branch-and-bound replaces exact optimization by a sound
under-approximation during search.

\begin{definition}[Lower bound]
A \emph{lower bound} for a retained multe-class \(C\) is a value
\(L(C)\) satisfying \(L(C)\le \operatorname{opt}(C)\).
\end{definition}

A lower bound need not equal the optimum.
It represents information known before optimization has been completed.

For additive cost, lower bounds compose naturally.

\begin{definition}[Compositional lower bound]
Suppose an e-node \(n\) references child multe-classes
\((C_1,\ldots,C_k)\).  Its induced lower bound is

\[
L(n)
=
c(n)
+
\sum_i
L(C_i).
\]
\end{definition}

Because every realization rooted at \(n\) contains \(n\) together
with one realization from each child class, \(L(n)\) is itself a
lower bound on every realization rooted at \(n\).

Suppose a realization of the root has already been found whose cost
is \(U\).  This incumbent provides an upper bound on the optimum.
If \(L(n)\ge U\), then no realization rooted at \(n\) can improve
upon the incumbent.  Such an e-node may therefore be discarded
without affecting the final optimum.

\begin{theorem}[Branch-and-bound soundness]

Pruning every e-node whose lower bound cannot improve upon the current
incumbent preserves the optimum root realization.

\end{theorem}

\paragraph{Discussion.}

Branch-and-bound introduces a second monotone computation into the
framework.  The semantic abstraction represented by multe-classes
computes a monotone approximation to semantic information: distinctions
may be forgotten, but never reintroduced.  Lower bounds compute a
monotone approximation to optimization information; rather than
attempting to compute the optimum cost directly, search accumulates
progressively stronger lower bounds on that optimum.

Each refinement establishes a stronger entailment
\(L(C)\le \operatorname{opt}(C)\).
The bound evolves monotonically until it either proves that a branch
cannot improve the incumbent or converges to the optimum itself.
Viewed this way, branch-and-bound is not merely an optimization
technique but an instance of monotone knowledge accumulation.

Semantic abstraction and lower-bound propagation therefore play
analogous roles: one computes increasingly precise information about
which realizations remain semantically distinguishable, the other about
which realizations remain competitively optimal.  Both are monotone
computations toward a fixed point.
% Exact minimum costs need not be known before pruning begins.
%
% Define a lower bound L(C) satisfying
%
%     L(C) <= opt(C).
%
% For an e-node n with references C_1,...,C_k, define its induced
% lower bound
%
%     L(n)
%       =
%     c(n) + sum_i L(C_i).
%
% Prove that L(n) is a lower bound on every realization rooted at n.
%
% Given an incumbent realization of cost U, n cannot improve the
% incumbent whenever
%
%     L(n) >= U
%
% under the chosen strict/non-strict tie policy.
%
% Distinguish two correctness principles:
%
% - Bellman compatibility justifies replacing a child by its best
%   representative;
%
% - bound soundness justifies discarding an alternative before its
%   exact optimum is known.
%
% State branch-and-bound correctness:
% pruning an e-node whose lower bound cannot beat the incumbent does
% not change the minimum achievable root cost.


%%%%%%%%%%%%%%%%%%%%%%%%%%%%%%%%%%%%%%%%%%%%%%%%%%%%%%%%%%%%%%%
\section{Examples}
\label{sec:examples}
%%%%%%%%%%%%%%%%%%%%%%%%%%%%%%%%%%%%%%%%%%%%%%%%%%%%%%%%%%%%%%%

The preceding sections introduced semantic schemas, multe-graphs,
requirements, canonicalization, and Bellman compatibility in the
abstract.  We now instantiate this framework in two familiar settings.
The first example shows that the framework recovers Selinger's
classical dynamic-programming optimizer; the second shows that it
strictly generalizes ordinary equality saturation by allowing several
semantic interfaces for the same equivalence class.

%%%%%%%%%%%%%%%%%%%%%%%%%%%%%%%%%%%%%%%%%%%%%%%%%%%%%%%%%%%%%%%
\subsection{Recovering Selinger's Optimizer}
\label{sec:selinger-example}
%%%%%%%%%%%%%%%%%%%%%%%%%%%%%%%%%%%%%%%%%%%%%%%%%%%%%%%%%%%%%%%

Consider optimization of a relational query.  Its semantic schema is
immediate.

\begin{center}
\begin{tabular}{ll}
Goal & logical relational expression \\
E-node & physical operator implementation \\
Composition & physical operator tree \\
\end{tabular}
\end{center}

A physical operator does not reference particular child plans; instead
it specifies the semantic properties required of its inputs.  A merge
join, for example, references child multe-classes providing compatible
sort orders, whereas a hash join references only unordered input
classes.  The refinement hierarchy records exactly the semantic
distinctions needed by these physical operators---in the classical
Selinger optimizer, the interesting orders.

A query becomes a problem instance by imposing requirements on the
root multe-class.
For example,

\[
\texttt{SELECT * FROM R JOIN S}
\]

requires any realization of the root goal, whereas

\[
\texttt{SELECT * FROM R JOIN S ORDER BY R.a}
\]

requires a root realization belonging to the
\([R\Join S]_{\mathsf{sorted}(a)}\)
multe-class.  Instantiation propagates this requirement downward
through the multe-graph, and only multe-classes that may contribute to
a satisfying realization remain live.

Canonicalization then chooses one representative for each surviving
multe-class.  Using additive plan cost and minimum-cost
canonicalization yields the lowest-cost implementation satisfying the
required semantic properties.

Bellman compatibility follows immediately.  The child references define
independent subproblems indexed by their required semantic interfaces,
and operator cost is monotone in child costs, so one minimum-cost
representative per surviving multe-class suffices.  The instantiated
multe-graph is therefore precisely the memoization structure maintained
by Selinger's dynamic programming algorithm.

Viewed in this way, interesting orders are not an auxiliary heuristic;
they are the semantic interfaces through which parent operators observe
their children.

%%%%%%%%%%%%%%%%%%%%%%%%%%%%%%%%%%%%%%%%%%%%%%%%%%%%%%%%%%%%%%%
\subsection{Generalizing Equality Saturation}
\label{sec:egraphs-example}
%%%%%%%%%%%%%%%%%%%%%%%%%%%%%%%%%%%%%%%%%%%%%%%%%%%%%%%%%%%%%%%

Equality saturation begins from a different semantic representation.
Each e-class represents all expressions known to be equivalent under
the rewrite system, and extraction chooses one representative
expression, typically by minimum cost.

In our terminology, an ordinary e-graph associates one semantic
quotient with each goal.  Every equivalent expression belongs to the
same e-class, every parent references that same class, and the semantic
interface is therefore fixed.

A multe-graph relaxes this restriction.  Instead of one semantic
interface, each goal may possess a hierarchy of semantic interfaces
represented by nested multe-classes, so different parent e-nodes may
observe different semantic distinctions among equivalent expressions.
One parent may distinguish expressions according to whether they
produce sorted output, for example, while another ignores output order
entirely; both observations concern the same logical result but require
different semantic interfaces.

Extraction proceeds exactly as before.  Requirements first instantiate
the multe-graph by removing inadmissible alternatives and unreachable
semantic distinctions, and canonicalization then chooses one
representative for each surviving multe-class.  Ordinary e-graph
extraction is recovered as the special case in which every goal
possesses only a single represented semantic quotient.

Multe-graphs therefore strictly generalize ordinary e-graphs,
preserving the fundamental equality-saturation model while allowing
multiple semantic interfaces to coexist for the same logical
equivalence class.

\paragraph{Discussion.}

The two examples illustrate complementary viewpoints.  Selinger's
optimizer begins with several semantic interfaces and builds a dynamic
program over them, whereas equality saturation begins with a single
semantic interface.

A multe-graph unifies these perspectives.  The refinement hierarchy
determines which semantic distinctions are represented, and the child
references determine the interfaces through which subproblems interact.
Bellman compatibility then guarantees that these interfaces are
sufficient for local optimization to compose into a globally optimal
solution.

From this perspective, dynamic programming is optimization over
semantic interfaces, and different optimization algorithms differ
primarily in how they choose those interfaces.

%%%%%%%%%%%%%%%%%%%%%%%%%%%%%%%%%%%%%%%%%%%%%%%%%%%%%%%%%%%%%%%
\section{Theory}
\label{sec:theory}
%%%%%%%%%%%%%%%%%%%%%%%%%%%%%%%%%%%%%%%%%%%%%%%%%%%%%%%%%%%%%%%

The preceding development separated compositional optimization into
five layers:

\begin{enumerate}

\item
semantic representation;

\item
problem instantiation;

\item
canonicalization;

\item
optimization; and

\item
search.

\end{enumerate}

This decomposition was introduced for conceptual clarity; the purpose
of the present section is to justify it mathematically.  Rather than
treating dynamic programming as a primitive algorithmic principle, we
derive it from properties of the underlying semantic representation.

The resulting theory has three themes.  First, we show that
multe-graphs faithfully represent the optimization problem.  Second, we
show that Bellman's Principle is a theorem about semantic interfaces
rather than an independent assumption.  Finally, we show that semantic
representation, optimization, and search have distinct correctness
obligations and therefore admit independent generalization.

%%%%%%%%%%%%%%%%%%%%%%%%%%%%%%%%%%%%%%%%%%%%%%%%%%%%%%%%%%%%%%%
\subsection{Representation}
\label{sec:representation-theory}
%%%%%%%%%%%%%%%%%%%%%%%%%%%%%%%%%%%%%%%%%%%%%%%%%%%%%%%%%%%%%%%

The first question is whether a multe-graph faithfully represents the
underlying optimization problem.  Instantiation specializes a semantic
schema to one problem instance by propagating requirements and removing
semantic distinctions that become unobservable; this transformation
should preserve exactly the realizations satisfying the specified
requirements.

\begin{theorem}[Representation correctness]
Extraction from the instantiated multe-graph is equivalent to
extraction from the original semantic schema restricted to realizations
satisfying the problem requirements.
\end{theorem}

Correctness alone is not sufficient.  Instantiation should also avoid
preserving distinctions that no satisfying realization can observe.

\begin{theorem}[Minimality]
Within the represented refinement hierarchy, the instantiated
multe-graph is the coarsest semantic quotient preserving every
distinction observable from the demanded root.
\end{theorem}

Together, these theorems justify the instantiated multe-graph as the
semantic representation of the optimization problem.

\begin{remark}[Future direction: minimal semantic interfaces]
The preceding theorem is global: it identifies the coarsest surviving
representation for an entire problem instance.  A complementary local
result appears possible.  Each child reference of a parent e-node is
intended to represent the coarsest semantic abstraction through which
that parent can correctly compose with its child.  Formalizing this
minimal-interface property would characterize multe-graphs as minimal
both globally---in the interfaces retained after instantiation---and
locally---in the information carried by each remaining interface.
\end{remark}

%%%%%%%%%%%%%%%%%%%%%%%%%%%%%%%%%%%%%%%%%%%%%%%%%%%%%%%%%%%%%%%
\subsection{Optimization}
\label{sec:optimization-theory}
%%%%%%%%%%%%%%%%%%%%%%%%%%%%%%%%%%%%%%%%%%%%%%%%%%%%%%%%%%%%%%%

Having established the semantic representation, we next ask when local
optimization is sufficient.  The key observation is that a parent
interacts with each child only through the referenced semantic
interface.  The semantic question is therefore whether replacing one
realization by another from the same referenced multe-class preserves
composition, while the optimization question is whether such a
replacement can worsen the objective.  These are independent concerns.

\begin{lemma}[Semantic substitution]
Replacing a child realization by another realization belonging to the
same referenced multe-class preserves semantic admissibility.
\end{lemma}

Semantic substitution is a property of the multe-graph alone,
independent of costs, objectives, and search.  Optimization enters only
through Bellman compatibility.

\begin{theorem}[Bellman's Principle]
Suppose

\begin{enumerate}

\item
referenced multe-classes are semantically substitutable;

\item
child subproblems are compositionally independent (the
Cartesian-fiber condition); and

\item
the optimization objective is monotone under semantic substitution.

\end{enumerate}

Then every optimal realization is assembled from optimal realizations
of its referenced child multe-classes.

Consequently, one optimal representative per retained multe-class
suffices for global optimization.
\end{theorem}

Bellman's Principle is therefore not assumed.  It follows from the
interaction between semantic interfaces and the optimization objective.

Branch-and-bound exploits the same interfaces in a different way.
Rather than computing optimum values directly, it accumulates monotone
lower bounds on those values.

\begin{theorem}[Branch-and-bound soundness]
Pruning an e-node whose lower bound cannot improve upon the current
incumbent preserves the optimum realization.
\end{theorem}

Unlike Bellman's Principle, this theorem depends only upon the
soundness of the lower bounds.

%%%%%%%%%%%%%%%%%%%%%%%%%%%%%%%%%%%%%%%%%%%%%%%%%%%%%%%%%%%%%%%
\subsection{Modularity}
\label{sec:modularity-theory}
%%%%%%%%%%%%%%%%%%%%%%%%%%%%%%%%%%%%%%%%%%%%%%%%%%%%%%%%%%%%%%%

The preceding results reveal that the framework separates several
independent correctness obligations.

\begin{center}
\begin{tabular}{ll}
Representation &
semantic interfaces preserve observable behavior
\\[1ex]

Instantiation &
requirements preserve satisfying realizations
\\[1ex]

Canonicalization &
one representative is selected per surviving interface
\\[1ex]

Optimization &
preferences compose across semantic interfaces
\\[1ex]

Search &
algorithms preserve the optimum
\end{tabular}
\end{center}

No theorem requires these layers to be developed simultaneously.
Changing the optimization objective does not alter semantic
correctness; changing the search algorithm does not alter Bellman
compatibility; and changing the refinement hierarchy changes the
represented semantic interfaces but not the meaning of
canonicalization.  Each layer therefore admits independent refinement.

\paragraph{Discussion.}

Traditional presentations of dynamic programming intertwine semantic
representation, optimization, and search.  The present development
separates these concerns: semantic interfaces determine what
information must flow between subproblems, optimization determines which
representatives of those interfaces are preferred, and search determines
how those preferences are computed.

Bellman's Principle is therefore not the foundation of compositional
optimization; rather, it is the theorem that semantic interfaces are
sufficient for independent optimization.  Viewed this way, dynamic
programming is optimization over semantic interfaces.

%%%%%%%%%%%%%%%%%%%%%%%%%%%%%%%%%%%%%%%%%%%%%%%%%%%%%%%%%%%%%%%
\section{Generalizations}
\label{sec:generalizations}
%%%%%%%%%%%%%%%%%%%%%%%%%%%%%%%%%%%%%%%%%%%%%%%%%%%%%%%%%%%%%%%

The preceding development deliberately adopted the simplest setting in
which the central ideas can be stated clearly, assuming tree-shaped
refinement hierarchies, root requirements, Boolean local predicates,
additive scalar cost, and finite acyclic multe-graphs.  None of these
assumptions is fundamental.  Once semantic representation,
instantiation, canonicalization, optimization, and search have been
separated, each may be generalized independently.  This section
sketches several directions.

%%%%%%%%%%%%%%%%%%%%%%%%%%%%%%%%%%%%%%%%%%%%%%%%%%%%%%%%%%%%%%%
\subsection{Beyond Tree Hierarchies}
%%%%%%%%%%%%%%%%%%%%%%%%%%%%%%%%%%%%%%%%%%%%%%%%%%%%%%%%%%%%%%%

Throughout this document, refinement hierarchies were represented as
trees.  This assumption simplifies presentation but is not essential.
More generally, a goal may admit many compatible semantic quotients
whose refinement order forms an arbitrary finite meet-semilattice or
partial order rather than a tree, so that different optimization
objectives may induce different quotient decompositions over the same
semantic space.

The local definitions of multe-classes and child references remain
unchanged; only the organization of the refinement hierarchy changes.
The representation and minimality theorems should therefore extend to
arbitrary compatible quotient families.

%%%%%%%%%%%%%%%%%%%%%%%%%%%%%%%%%%%%%%%%%%%%%%%%%%%%%%%%%%%%%%%
\subsection{Beyond Root Requirements}
%%%%%%%%%%%%%%%%%%%%%%%%%%%%%%%%%%%%%%%%%%%%%%%%%%%%%%%%%%%%%%%

The presentation treated problem requirements as originating at the
root, but in practice requirements naturally arise throughout the
graph.  Examples include

\begin{itemize}

\item memory budgets for individual operators,

\item physical constraints on particular subplans,

\item user hints,

\item cached intermediate results,

\item preferred rewrite strategies.

\end{itemize}

The semantic framework already accommodates this generalization.  A
requirement is simply a predicate registered with a multe-class or
Local predicates depend only upon the intrinsic properties of one
e-node.  Many practically important requirements instead constrain
aggregate properties of complete realizations.  Examples include
Local predicates depend only upon the intrinsic properties of one
e-node.

Many practically important requirements instead constrain aggregate
properties of complete realizations.

Examples include

\begin{itemize}

\item total execution cost,

\item peak memory consumption,

\item latency,

\item energy usage.

\end{itemize}

Unlike local predicates, aggregate requirements propagate through the
composition equations defining the corresponding attribute, and exact
propagation generally produces relations among child demands rather
than independent child requirements.  Approximate propagation, lower
The running development specialized optimization to additive scalar
cost, but the essential property used by Bellman's Principle is weaker.
Each e-node may instead define an operator-specific aggregation
function \(F_n:A^k\rightarrow A\)
over an arbitrary attribute domain \(A\).  Bellman compatibility then
becomes a property of these aggregation functions rather than of
addition itself, with monotonicity under semantic substitution
replacing additivity as the fundamental optimization requirement.
Uniform algebraic structures such as semirings and dioids provide
important special cases, but are not required by the semantic
framework.rather than of addition itself.

Monotonicity under semantic substitution replaces additivity as the
fundamental optimization requirement.

Uniform algebraic structures such as semirings and dioids provide
important special cases, but are not required by the semantic
framework.

%%%%%%%%%%%%%%%%%%%%%%%%%%%%%%%%%%%%%%%%%%%%%%%%%%%%%%%%%%%%%%%
\subsection{Beyond Acyclicity}
%%%%%%%%%%%%%%%%%%%%%%%%%%%%%%%%%%%%%%%%%%%%%%%%%%%%%%%%%%%%%%%

Finite acyclic multe-graphs permit inductive definitions of
realizability, instantiation, and extraction.  Recursive optimization
problems instead require cyclic semantic schemas, and the inductive
definitions then become least-fixed-point definitions, with requirement
propagation, optimization, and lower-bound refinement all becoming
monotone operators over suitable lattices.  The present framework
therefore extends naturally toward recursive queries, recursive rewrite
systems, and general fixed-point semantics.

%%%%%%%%%%%%%%%%%%%%%%%%%%%%%%%%%%%%%%%%%%%%%%%%%%%%%%%%%%%%%%%
\subsection{Beyond Scalar Optimization}
%%%%%%%%%%%%%%%%%%%%%%%%%%%%%%%%%%%%%%%%%%%%%%%%%%%%%%%%%%%%%%%

Scalar optimization assumes that every realization receives a value in
a totally ordered domain, yet many optimization problems are inherently
multi-objective.  Rather than selecting a single scalar optimum,
canonicalization may operate over a lifted outcome domain whose elements
are Pareto frontiers, with the frontier itself serving as one canonical
outcome in that domain.  Pareto optimization thus changes the outcome
space rather than the definition of canonicalization, and Bellman
compatibility again becomes a question about monotonicity in the lifted
domain.

%%%%%%%%%%%%%%%%%%%%%%%%%%%%%%%%%%%%%%%%%%%%%%%%%%%%%%%%%%%%%%%
\subsection{Beyond Optimization}
%%%%%%%%%%%%%%%%%%%%%%%%%%%%%%%%%%%%%%%%%%%%%%%%%%%%%%%%%%%%%%%

The preceding sections implicitly assumed that canonicalization is
entailed by the optimization objective.  Scalar optimization chooses a
unique optimum, and Pareto optimization a unique Pareto frontier in a
lifted domain; in both cases the canonical outcome is determined by
semantic entailment.

Optimization is therefore only one canonicalization discipline.  More
generally, a semantic interface may admit several incomparable canonical
outcomes even after every available optimization criterion has been
applied.  Choosing among those outcomes is no longer optimization but a
semantic commitment.

This observation motivates Determination Theory.  The semantic
representation developed in the present document is independent of the
canonicalization policy applied to it, of which optimization is only
one.  Determination studies the cases in which semantic entailment alone
is insufficient to select a unique representative.

%%%%%%%%%%%%%%%%%%%%%%%%%%%%%%%%%%%%%%%%%%%%%%%%%%%%%%%%%%%%%%%
\section{Payoffs}
\label{sec:payoffs}
%%%%%%%%%%%%%%%%%%%%%%%%%%%%%%%%%%%%%%%%%%%%%%%%%%%%%%%%%%%%%%%

The development above is deliberately spare: it fixes a vocabulary
and separates concerns.  This final section collects the results
that the framework exists to support.  \emph{In an article prepared
for submission, the material of this section would be the prominent
results}; the preceding sections would be compressed into the
definitions needed to state them.  We record each payoff here at
internal-note fidelity, marking claims as theorems, conjectures, or
open problems according to our current confidence.

The intended spine of the outbound article is:

\begin{enumerate}
\item \emph{semantic interfaces} --- the framework of the preceding
      sections;
\item \emph{their derivation} --- continuation equivalence and a
      Myhill--Nerode theorem for dynamic programming
      (Section~\ref{sec:continuation}); and
\item \emph{their price} --- width and dimension as intrinsic
      complexity measures of compositional optimization problems
      (Section~\ref{sec:width});
\end{enumerate}

with demand transformation (Section~\ref{sec:demand}), local search
(Section~\ref{sec:terrain}), and rewriting-system completion
(Section~\ref{sec:completion}) as applications.  A closing
dictionary (Section~\ref{sec:types}) translates the framework into
the vocabulary of programming languages, where several of our
central notions carry standard names.

Throughout, one sentence serves as motif:

\begin{quote}
\emph{A semantic interface is exactly the information about a
subproblem that every future composition may need.}
\end{quote}

Everything in this document---Selinger, Bellman, e-graphs,
branch-and-bound, demand propagation, local search, and the
determination boundary---is a consequence of that one idea.

%%%%%%%%%%%%%%%%%%%%%%%%%%%%%%%%%%%%%%%%%%%%%%%%%%%%%%%%%%%%%%%
\subsection{Deriving Semantic Interfaces}
\label{sec:continuation}
%%%%%%%%%%%%%%%%%%%%%%%%%%%%%%%%%%%%%%%%%%%%%%%%%%%%%%%%%%%%%%%

The preceding development takes the refinement hierarchies as given:
multe-classes are supplied with the schema, and the theorems of
Section~\ref{sec:theory} assume them.  The deepest question the
framework raises is where they come from.  This subsection derives
them.  In the submission version we expect the results sketched
here---a Myhill--Nerode theorem for dynamic programming---to be the
central theoretical contribution.

The role of ``the future of a realization'' is played by
\emph{contexts}: realizations of the root goal with a hole into
which a realization of $x$ may be plugged.  For each goal $x$, let
$\mathcal T(x)$ denote the set of complete concrete realizations of
$x$ (Section~\ref{sec:extraction}).

\begin{definition}[Context]
Let $x\in X$ be a goal.  A \emph{context} for $x$ is a complete
realization of the root goal $x_\star$ with a single hole of sort
$x$: formally, an element of $\mathcal T(x_\star)$ in the schema
extended with a fresh atomic e-node
$[\,\cdot\,]_x\in\mathsf N(x)$ that is used exactly once.
Write $\mathsf{Ctx}(x)$ for the set of contexts for $x$, and $C[t]$
for the term obtained by plugging $t\in\mathcal T(x)$ into the hole.
\end{definition}

\begin{definition}[Acceptance]
A context $C\in\mathsf{Ctx}(x)$ \emph{accepts} $t\in\mathcal T(x)$
when $C[t]$ is an admissible realization of $x_\star$: every
composition step of $C[t]$ lies in its composition relation.
When a problem instance is fixed, we additionally require that the
root e-node of $C[t]$ satisfy the root requirement $\rho_\star$.
\end{definition}

\begin{definition}[Continuation equivalence]
\label{def:continuation}
For $s,t\in\mathcal T(x)$, define $s\equiv_x t$ when

\begin{enumerate}
\item every context accepts $s$ exactly when it accepts $t$; and
\item for every accepting context $C$,
\[
\operatorname{cost}(C[s])-\operatorname{cost}(s)
=
\operatorname{cost}(C[t])-\operatorname{cost}(t).
\]
\end{enumerate}

The quotient
\[
K_x=\mathcal T(x)/{\equiv_x}
\]
is the \emph{continuation quotient} at $x$.
\end{definition}

Two realizations are continuation-equivalent when swapping one for
the other inside any complete realization changes nothing that the
rest of the realization can observe: no context can tell them apart,
so \emph{no optimizer ever needs to distinguish them}.  Readers who
recall automata theory will recognize a Myhill--Nerode construction:
it quotients realizations by their future behavior, just as
Myhill--Nerode quotients string prefixes by their acceptance
futures.

\begin{remark}[Additive collapse]
Under the additive cost model of
Section~\ref{sec:compositional-cost}, the incremental-cost condition
is vacuous: $\operatorname{cost}(C[s])-\operatorname{cost}(s)$ is
the sum of the local costs of the e-nodes of $C$ itself, independent
of $s$.  Continuation equivalence is then a purely semantic
relation---two realizations are equivalent exactly when the same
contexts accept them---and cost enters only through dominance among
inequivalent classes.  This matches practice: interesting orders are
semantic properties, not costs.  Under the general aggregation
functions of Section~\ref{sec:generalizations}, the cost condition
has content: a realization may be accepted everywhere yet contribute
differently to different parents.
\end{remark}

Continuation equivalence gives a yardstick for hierarchies.

\begin{definition}[Sufficient and exact hierarchies]
A refinement hierarchy for a goal $x$ is \emph{Bellman-sufficient}
when any two realizations represented by the same leaf multe-class
are continuation-equivalent, and \emph{exact} when, in addition, two
realizations are represented by the same leaf class if and only if
they are continuation-equivalent.
\end{definition}

Selinger's interesting orders~\cite{selinger1979access} are a
Bellman-sufficient hierarchy for the System~R cost model; whether
they are exact depends on the operator inventory.

\begin{conjecture}[Minimality]
\label{conj:minimality}
The continuation quotient is the coarsest Bellman-sufficient
interface: every Bellman-sufficient hierarchy refines $K_x$, the
factoring is unique, and the family of quotient maps commutes with
child references and with instantiation.  Consequently, every
Bellman-compatible multe-graph factors through the continuation
quotients, which are initial among Bellman-valid semantic
representations.
\end{conjecture}

The Minimality theorem of Section~\ref{sec:representation-theory} is
the instance-level shadow of this conjecture: instantiation discards
exactly the distinctions that no accepting context of the instance
can observe, which is coarsening toward the instance-relative
continuation quotient.  The conjecture strengthens that statement
from the represented hierarchy to all hierarchies.

\begin{conjecture}[Universality]
\label{conj:universality}
Let a concrete semantic schema with costs be given, with no
refinement hierarchy.  Then the continuation quotients themselves
support a multe-graph: composition descends to the quotients, the
descended fibers admit Cartesian presentations (after the usual
splitting of e-nodes), and the induced objective is monotone under
substitution within classes.  Consequently, a compositional
optimization problem admits an exact interface-indexed dynamic
program if and only if its continuation quotients are uniformly
small.
\end{conjecture}

\begin{remark}[A Myhill--Nerode theorem for dynamic programming]
Together, Conjectures~\ref{conj:minimality}
and~\ref{conj:universality} assert: \emph{Bellman's principle holds
exactly when future distinguishability has a finite compositional
representation, and the continuation quotient is its minimal
automaton.}  Selinger's interesting orders, Cascades' property
requests~\cite{graefe1995cascades}, and e-graph extraction
frontiers~\cite{willsey2021egg} become different coordinate systems
on the same canonical object; interesting orders, in particular, are
the Nerode classes of physical plans.  This is the strongest sense
in which the multe-graph is derived rather than designed.
\end{remark}

\begin{remark}[Two presentations of factorization]
\label{rem:presentations}
The Bellman condition has two equivalent presentations, and the
universality proof should fix one and prove the equivalence once.
The \emph{relational} presentation is the one used in this document:
composition relations with the Cartesian-fiber condition of
Section~\ref{sec:multe-graph-definition}, where a fiber that is not
one product is split across several e-nodes.  The \emph{functional}
presentation labels each realization with its interface and equips
each operator with a property-transfer map computing the parent
label from the child labels; partial operators are accommodated by
splitting them per admissible combination of child labels.  The two
splittings are the same maneuver.  The relational form is more
general in appearance; the functional form makes the congruence
argument cleaner.
\end{remark}

\begin{remark}[Deferred commitment is a free completion]
\label{rem:free}
The Pareto lift of Section~\ref{sec:generalizations} and the
frontier-shaped memo structures of practical optimizers share a
canonicity property worth recording, because it shows they are
forced rather than chosen.  Let $(O,\preceq)$ be a finite poset of
properties.  The domain of joint demands is $\mathsf{Down}(O)$, the
finite downsets of $O$ ordered by inclusion, and this is the
\emph{free} join-semilattice with bottom on $O$: every monotone map
from $O$ into a join-semilattice extends uniquely to a
join-preserving map on $\mathsf{Down}(O)$.  The completion therefore
adds exactly one capability to the property domain---holding several
incomparable alternatives open at once---and nothing else.  Any
smaller summary must conflate alternatives that some composition
treats differently.  In particular, collapsing all incomparable sets
to a single formal ``undetermined'' element identifies $\{a,b\}$,
$\{a,c\}$, and $\{a,b,c\}$, so transfer maps can no longer act
differently on them and factorization fails for any operator that
consumes one alternative but not another.  The downset completion is
the least completion compatible with deferred commitment.
\end{remark}

\TODO{Nested antichains (``ears on ears'').  In a join-semilattice,
nested joins flatten: $a\sqcup(b\sqcup c)=a\sqcup b\sqcup c$, so
nested antichains such as $\{\{c,d\},b\}$ denote nothing beyond
their flattening $\{b,c,d\}$ in a demand domain.  Earlier discussion
suggested that nesting records the dependency structure among future
commitments (staged or conditional choices), which---if it must be
retained---belongs in the composition structure or in a finer
multe-class tree rather than in the demand domain.  Decide: is the
flat completion always enough, with staging carried by the goals and
composition relations, or do we need an explicit AND/OR term algebra
over joins and compositions?  This may also require a fix in the
determination papers.  Resolve before the universality proof
attempt.}

%%%%%%%%%%%%%%%%%%%%%%%%%%%%%%%%%%%%%%%%%%%%%%%%%%%%%%%%%%%%%%%
\subsection{The Price of Bellman: Width and Dimension}
\label{sec:width}
%%%%%%%%%%%%%%%%%%%%%%%%%%%%%%%%%%%%%%%%%%%%%%%%%%%%%%%%%%%%%%%

This subsection is the second intended headline of the outbound
article.  Its premise converts the framework from a qualitative
condition into a quantitative theory: Bellman's principle is never
unavailable, only expensive, and the expense is an intrinsic
complexity measure of the problem.

The claim of availability is not a conjecture; it is a nearly
trivial proposition, and its triviality is the point.

\begin{proposition}[Sufficiency is always achievable]
\label{prop:trivial}
Every concrete semantic schema with costs admits a
Bellman-compatible multe-graph presentation.
\end{proposition}

\begin{proof}[Proof sketch]
Take the discrete presentation: unfold the schema until every
realization is represented by its own chain of e-nodes, and let
every multe-class be a singleton.  Each composition fiber is then a
union of $1\times\cdots\times 1$ products, one e-node per tuple, so
the Cartesian-fiber condition holds trivially; substitution within a
singleton class is the identity, so monotonicity is vacuous.
\end{proof}

This lift is useless in practice: the representation is as large as
the solution space itself.  Its lesson stands nonetheless.
Sufficiency is automatic, so the substantive question is
quantitative: \emph{how much must be retained?}  The continuation
quotient supplies the yardstick.

\begin{remark}[Choosing the decomposition]
One honest hedge applies to everything that follows: the framework
takes the decomposition---the goals and composition signatures---as
given, and a full theory of choosing good decompositions remains
open.  The hedge is only partial.  Completion as preprocessing
(Section~\ref{sec:completion}) rewrites a given decomposition into
one that is as monotone as the objective allows: the orientable
portion is discharged by normalization, and multe-class structure is
confined to the residue where the order genuinely cannot decide.
\end{remark}

\begin{definition}[Bellman width]
\label{def:width}
The \emph{width} of a goal $x$ is $w(x)=|K_x|$, the number of
continuation classes at $x$.  The width of a problem is
$\max_{x\in X}w(x)$.
\end{definition}

\begin{definition}[Bellman dimension]
\label{def:dimension}
A \emph{property vocabulary} is a finite poset $(O,\preceq)$
together with a labeling assigning each realization the downset of
properties it provides, such that labels factor through composition
in the sense of the functional presentation
(Remark~\ref{rem:presentations}).  The \emph{dimension} of a problem
is the least $|O|$ over vocabularies whose labeling is
Bellman-sufficient (equal labels imply continuation equivalence);
the \emph{Boolean dimension} restricts to discretely ordered $O$.
\end{definition}

\begin{proposition}[Basic bounds; proofs to be recorded]
\label{prop:bounds}
\begin{enumerate}
\item Every Bellman-sufficient hierarchy at $x$ has at least $w(x)$
  leaf classes; hence $w(x)$ lower-bounds the size of every exact
  memo structure at $x$.
\item Boolean labels at $x$ are subsets of the properties relevant
  at $x$, so the Boolean dimension is at least
  $\max_x\lceil\log_2 w(x)\rceil$.
\item Extraction runs in time proportional to
  $\sum_{n}\prod_{i} f(C_i)$ over non-atomic e-nodes $n$ with child
  references $(C_1,\ldots,C_k)$, where $f(C_i)$ is the size of the
  retained frontier at $C_i$.
\end{enumerate}
Consequently, dynamic programming is efficient exactly when widths
are small: polynomial width gives polynomial memo structures, and
exponential width forbids them.
\end{proposition}

\begin{example}[Held--Karp width]
\label{ex:heldkarp}
The traveling salesperson problem fits the framework with goals
``simple paths from a fixed start visiting exactly $k$ cities,''
e-nodes extending a path by one city, and complete tours as the root
goal.  For generic edge weights, two partial paths are
continuation-equivalent if and only if they visit the same city set
and end at the same city, so
\[
w(x_k)=\binom{n}{k}\,k :
\]
the width is exponential in $n$.  The Held--Karp dynamic
program~\cite{heldkarp1962} memoizes exactly these classes, so it is
the canonical lift---and its $2^n$ table is not an artifact of a
weak state choice but the price the problem demands, at least among
interface-indexed dynamic programs.  Selinger's problem, by
contrast, has width bounded by the number of realizable order
bundles, which is small for typical queries.  Characterizing which
problems have polynomial width is, we believe, the right abstract
form of the question ``when does dynamic programming work?''
\end{example}

\begin{remark}[Rectangle covers and communication complexity]
\label{rem:rectangles}
The relational presentation supplies a second size measure and,
with it, lower-bound technology.  A non-atomic e-node presents one
Cartesian product of child classes, so a concrete composition fiber
must be covered by such products: the minimum number of e-nodes
presenting a fiber is its \emph{rectangle-cover number} with respect
to the child quotients---in other words, a nondeterministic
communication complexity.  Width counts classes per goal; covers
count e-nodes per composition.  Communication-complexity lower
bounds should therefore transfer directly to multe-graph size,
giving problems for which \emph{any} Bellman-compatible
representation is provably large.  We regard this as the most
promising route to lower bounds for the open problems below.
\end{remark}

\begin{openproblem}[Compositionality gap]
\label{open:gap}
The bound of Proposition~\ref{prop:bounds}(2) is
information-theoretic: $\lceil\log_2 w\rceil$ properties suffice to
\emph{name} the continuation classes, but a vocabulary must also
factor through every composition, which may force more properties.
Determine how large the gap between the compositional dimension and
$\lceil\log_2 w\rceil$ can be, and identify structural conditions
under which the information-theoretic bound is achieved.
\end{openproblem}

\begin{openproblem}[Small-width problems]
\label{open:width}
Characterize the compositional problems of polynomial width.  Such a
characterization would delimit, in one vocabulary, the problems for
which Selinger-style dynamic programming, Cascades-style
memoization, and frontier-indexed e-graph extraction can be
efficient.
\end{openproblem}

%%%%%%%%%%%%%%%%%%%%%%%%%%%%%%%%%%%%%%%%%%%%%%%%%%%%%%%%%%%%%%%
\subsection{Demand Transformation}
\label{sec:demand}
%%%%%%%%%%%%%%%%%%%%%%%%%%%%%%%%%%%%%%%%%%%%%%%%%%%%%%%%%%%%%%%

This subsection and the next are payoffs of a different kind: in the
outbound article they would form the applications section, showing
that semantic interfaces organize search architectures as well as
correctness.  The observation here is that \emph{instantiation is
demand propagation}, and demand propagation has a classical theory.

Recall the two-directional propagation of
Section~\ref{sec:propagation}: liveness propagates demands downward
from the accepted root e-nodes, realizability computes supply upward
from the atomic e-nodes, and the instantiated multe-graph is their
intersection.  In deductive databases this construction has a name.
Magic-sets rewriting~\cite{bancilhon1986magic} and its refinement,
demand transformation with subsumptive
tabling~\cite{tekle2011subsumptive}, specialize a bottom-up program
so that it computes only facts demanded by a calling context.
There, an \emph{adornment} pairs a predicate with a demand
descriptor exactly as an instantiated demand pairs a multe-class
with a requirement.  One notable difference: magic sets specializes
on \emph{values} (which tuples) to prune irrelevant computation,
while we specialize on \emph{semantic properties} (how outputs are
shaped) to restore Bellman's principle---the same rewrite, deployed
for relevance in one case and sufficiency in the other.

We see three uses for the correspondence.

First, it derives the demands.  Define the \emph{demand closure} as
the least set of (class, requirement) pairs containing the root
demand and closed downward through the child references of the
e-nodes surviving each demand.  Restricting extraction to the
closure should preserve the root value, and the union of demands
registered at each goal is a candidate constructive answer to which
distinctions matter there---an input to the universality proof
(Conjecture~\ref{conj:universality}).

Second, it names a duality we have circled repeatedly:
\emph{Cascades is the magic-sets rewriting of Selinger}.  Top-down
optimizers propagate required physical properties from parent to
child with memoization; that is demand-transformed evaluation of the
bottom-up recurrence.  The slogan is implicit in prior work: the
Evita Raced metacompiler~\cite{condie2008evita} implemented both a
Selinger-style bottom-up optimizer and a Cascades-style top-down
optimizer as Overlog programs over a single declarative runtime.
What was there an engineering observation---the two optimizers are
small variations of one declarative program---becomes here a theorem
target: the top-down architecture should be \emph{derivable} from
the bottom-up one by a demand rewrite with a classical correctness
proof pattern.

Third, it explains an optimization.  Answering a demand from a
stronger cached supply is subsumption, and the result that
subsumptive tabling beats magic sets~\cite{tekle2011subsumptive}
maps onto our satisfaction relation: dominance-aware memoization is
exactly the optimization their analysis rewards.

\begin{remark}[Demand can overshoot supply]
The demand closure ranges over all child demands admitted by the
composition structure and may therefore demand classes whose
realizable restriction is empty.  The practically right object is
the intersection of demand, computed top-down, with realized supply,
computed bottom-up---which is exactly what instantiation already
computes, and what a Cascades-style memo achieves operationally.
Any minimality claim about derived demands must be stated relative
to this intersection, not to demand alone.
\end{remark}

\begin{openproblem}[Demand transformation of extraction]
\label{open:demand}
Formalize the demand closure, prove that restriction to it preserves
the root value, show that the induced relevance assignment is least
among antitone assignments supporting the root demand, and determine
whether the complexity guarantees of demand transformation and
subsumptive tabling transfer to interface-indexed dynamic programs.
\end{openproblem}

%%%%%%%%%%%%%%%%%%%%%%%%%%%%%%%%%%%%%%%%%%%%%%%%%%%%%%%%%%%%%%%
\subsection{Local Search and Cost Terrains}
\label{sec:terrain}
%%%%%%%%%%%%%%%%%%%%%%%%%%%%%%%%%%%%%%%%%%%%%%%%%%%%%%%%%%%%%%%

This payoff is the most speculative material in the document; in
the outbound article it would be a closing vision section, or a
separate paper.  Its role is to show that semantic interfaces
organize not only exact optimization but heuristic search over the
same problems.

A \emph{cost terrain} is a graph whose points are complete
realizations $\mathcal T(x_\star)$, with an edge between two points
when a single local move transforms one into the other, and with the
objective as elevation.  The key observation is that \emph{local
moves are substitutions}: every practical move---a swap, a rotation,
a 2-opt exchange---rewrites the current point $C[s]$ as $C[t]$ for
some context $C\in\mathsf{Ctx}(x)$ and replacement
$t\in\mathcal T(x)$.  Move sets restrict which goals $x$ and which
replacements $t$ are eligible, but every move has this shape.  In
the language of this document:

\begin{quote}
Every local-search operator is the replacement of one representative
of a semantic interface by another.  Bellman-style optimization and
local search use the same primitive operation.  One chooses globally
optimal substitutions through sufficient interfaces; the other
explores substitutions heuristically, one commitment at a time.
\end{quote}

\begin{lemma}[Substitution monotonicity; to be proved]
\label{lem:terrain-mono}
Under Bellman compatibility, replacing a piece $s$ by a
continuation-equivalent piece $t$ with
$\operatorname{cost}(t)\le\operatorname{cost}(s)$ never increases
the total cost: $\operatorname{cost}(C[t])\le\operatorname{cost}(C[s])$
for every accepting context $C$.  (Induct up the spine of $C$: at
each composition the child interfaces agree, so admissibility is
preserved, and the costs compare by monotonicity of the objective.)
Consequently, a move that lowers a piece's cost can raise the total
cost only by \emph{weakening the piece's interface}.
\end{lemma}

The lemma locates the source of ruggedness.  A greedy improver stuck
short of the optimum---despite having improving pieces
available---is stuck on an \emph{interface tradeoff}: a
cheaper-but-weaker piece against a costlier-but-stronger one.  These
are exactly the dominance-incomparable entries of a Pareto frontier;
the branching antichains of the cost--interface frontier are the
generators of the terrain's non-convexity.  If every goal's
realizations carried a single interface, piecewise improvement would
compose monotonically and the ridges would vanish.  In this precise
sense, \emph{ruggedness is width made visible}: the notoriously
rugged move terrains of the traveling salesperson problem and
Held--Karp's exponential table (Example~\ref{ex:heldkarp}) are two
symptoms of one underlying quantity.

One distinction keeps the claim honest.  A point can be locally
minimal for two different reasons: every available move loses on a
cost--interface tradeoff, or the winning substitution simply lies
outside the move set, as when a 2-opt local minimum is escaped by a
3-opt move.  The first reason is intrinsic to the problem and is
what the width analysis measures; the second is an artifact of
search design, the classical concern of neighborhood
exactness~\cite{papadimitriou1982combinatorial}, and our framework
does not speak to it.

\begin{remark}[Deferred commitment dissolves the terrain]
A local search occupies one terrain point at a time, so at every
interface tradeoff it makes an irrevocable commitment.  A
frontier-indexed dynamic program instead carries, per goal, the
entire cost--interface frontier: the discrete choice among
incomparable pieces is deferred until the context that discriminates
among them has been assembled.  This is an order-theoretic analogue
of convexification: convexification enlarges a continuous search
space until local optima disappear; the frontier lift enlarges the
memo state until Bellman's principle holds.  Equivalently, one may
keep the terrain---its points and moves unchanged---and lift the
\emph{cost}: replace the scalar objective by a vector objective
compared by dominance, from which the scalar is recovered by a fixed
monotone aggregation.  We call this a \emph{cost lift}; its size is
the number of cost dimensions.  Taking one coordinate per maximal
interface---in the spirit of Ioannidis and Kang's analysis of query
optimization terrains~\cite{ioannidis1991left}---the cost lift and
the semantic lift become the same construction in different
coordinates.
\end{remark}

\begin{remark}[Dimension as terrain smoothing]
\label{rem:terrain-dim}
The terrain gives the Bellman dimension
(Definition~\ref{def:dimension}) a reading in the geometry of the
cost surface itself, independent of any property vocabulary.  Call a
cost lift \emph{smoothing} when dominance-descent by substitution
never strands: from every suboptimal point, some move improves in
the dominance order.  By Lemma~\ref{lem:terrain-mono}, interface
tradeoffs are the only ridges, so smoothing requires exactly that
every genuine tradeoff remain visible as an incomparability of
lifted points rather than collapse into a scalar ridge.  Realizing a
prescribed incomparability structure by dominance in $d$ scalar
coordinates is precisely the classical \emph{order dimension} of
Dushnik and Miller~\cite{dushnikmiller1941}.  We therefore expect
the minimal smoothing dimension to be the order dimension of the
forced-preference structure on continuation classes.  The Boolean
dimension bounds it from above, since bundle labels embed the
frontier in a Boolean lattice whose order dimension is exactly
$|O|$; Sperner's theorem gives the counting form of the lower bound.
The compositionality gap (Open Problem~\ref{open:gap}) thereby
acquires a geometric face: the coordinates of a smoothing lift must
be features that substitution respects, propagating through every
composition, so the compositional dimension may exceed the order
dimension, and by how much is the gap.
\end{remark}

\begin{remark}[Basin count is not dimension]
It is tempting to conclude that dimension is a function of the
number of basins in the terrain, but the two quantities are
classically independent: posets of order dimension two already
contain arbitrarily large antichains, so two coordinates can smooth
a terrain with arbitrarily many mutually-ridged basins.  Moreover
the basin count varies with the move set and the particular costs,
which the order dimension does not.  The distinction is between
counting and structure: basin count asks \emph{how many} tradeoffs
the instance realizes, while dimension asks how the tradeoffs
\emph{interlock}---whether every incomparability can be explained by
the same two rankings disagreeing, or whether the incomparabilities
force many mutually independent rankings, as in the standard
examples of dimension theory~\cite{trotter1992}.
\end{remark}

\begin{openproblem}[Geometric dimension of a dynamic program]
\label{open:geometric-dim}
Make the smoothing picture precise.
\begin{enumerate}
\item \emph{Identify the relation.}  Order dimension realizes a
  given partial order; our constraints prescribe certain
  incomparabilities while leaving other comparabilities free.  The
  quantity at stake is the minimum dimension over all partial orders
  consistent with the constraints.  Define it, and determine whether
  it coincides with a standard notion.
\item \emph{Prove the sandwich.}  Conjecturally, the order dimension
  of the forced structure is at most the compositional (Bellman)
  dimension, which is at most $|O|$ for every sufficient vocabulary
  $O$, with the first inequality strict exactly when
  substitution-respecting coordinates cost more than unconstrained
  ones (Open Problem~\ref{open:gap}).
\item \emph{Settle the complexity.}  Deciding order dimension
  $\le 2$ is polynomial; $\le 3$ is
  NP-complete~\cite{yannakakis1982}.  Determine what this implies
  for computing or approximating the Bellman dimension, and identify
  tractable classes.
\end{enumerate}
\end{openproblem}

Several bodies of literature meet here, and each offers something
specific.  Dimension theory for finite posets~\cite{trotter1992}
supplies lower-bound technology (critical pairs, the standard
examples) that should transfer to Bellman dimension essentially
unchanged.  The complexity theory of local
search---PLS~\cite{johnson1988easy} and neighborhood
exactness~\cite{papadimitriou1982combinatorial}---addresses how hard
it is to \emph{find} local optima, where our width analysis
addresses \emph{why} they exist; whether small width places a
problem in a tractable fragment of PLS is, to our knowledge, open.
Finally, the empirical landscape-analysis tradition in query
optimization~\cite{ioannidis1991left} measured the shapes of real
cost surfaces; width and dimension are candidate structural
explanations for those measurements, and re-running such experiments
with frontier-indexed state would test the theory.

%%%%%%%%%%%%%%%%%%%%%%%%%%%%%%%%%%%%%%%%%%%%%%%%%%%%%%%%%%%%%%%
\subsection{Completion and Rewriting}
\label{sec:completion}
%%%%%%%%%%%%%%%%%%%%%%%%%%%%%%%%%%%%%%%%%%%%%%%%%%%%%%%%%%%%%%%

This payoff is addressed to readers arriving from term rewriting
rather than from query optimization.  In the outbound article it
would extend the width story to a third community---after databases
and e-graphs---and, as discussed below, it supplies a candidate
algorithmic skeleton for the universality proof.  The claim is that
Knuth--Bendix completion~\cite{knuthbendix1970,baadernipkow1998} is
not an analogy for this framework but a corner of it: the width-one
corner, degenerate in exactly the way that makes the rest of the
theory legible.

The core observation concerns reduction orders.  The defining
requirement on a reduction order---closure under contexts,
$s>t\Rightarrow C[s]>C[t]$---is verbatim the substitution
monotonicity of Lemma~\ref{lem:terrain-mono}, demanded for
\emph{every} context rather than derived for sufficient interfaces.

\begin{quote}
\emph{A reduction order is exactly a cost model for which Bellman
compatibility holds at the coarsest possible interface.}
\end{quote}

Under such an objective every class has width one: no two contexts
ever disagree about a pair of equivalent terms, so one canonical
representative per class suffices and no refinement hierarchy is
needed.  The Knuth--Bendix order itself makes the correspondence
nearly literal: symbol weights summed over a term, with
lexicographic tie-breaking, is an additive local cost
(Section~\ref{sec:compositional-cost}) constrained to be
context-monotone.

The rest of the correspondence is worth recording entry by entry,
because each entry suggests a theorem.

\begin{itemize}

\item
\emph{A convergent rewriting system is a compositionally coherent
canonicalization policy.}  Normal forms are closed under subterms,
so the canonical representative of a class is assembled from
canonical representatives of child classes: Bellman's principle for
the order.

\item
\emph{Termination is well-foundedness.}  The reduction order
certifies that induction is available, playing the role that
acyclicity plays in this document; the fixed-point generalization of
Section~\ref{sec:generalizations} corresponds to dropping
termination, which is what saturation does.

\item
\emph{Local confluence is the local coherence test.}  Confluence
guarantees that canonicalization is \emph{path}-independent; Bellman
compatibility guarantees that it is \emph{context}-independent.
Critical pairs are a finite local test for a global coherence
property, via Newman's lemma~\cite{newman1942}---structurally the
same move as testing the congruence condition e-node by e-node
rather than realization by realization.

\item
\emph{A non-orientable equation is a dominance-incomparable pair.}
When neither $s>t$ nor $t>s$ holds in all contexts, width greater
than one is surfacing.  Classical completion fails here; unfailing
completion~\cite{bachmair1989completion} keeps the equation and
orients per instance---an ad hoc, runtime version of orienting per
interface.  The multe-class machinery is the principled version:
split the class by the contexts that disagree, and keep one
canonical form per interface.

\item
\emph{Orientation decisions the order does not entail are
determinations.}  When completion requires the user to choose an
orientation, that is a semantic commitment beyond entailment---the
boundary drawn at the end of Section~\ref{sec:generalizations},
encountered in 1970 under a different name.  Rewriting modulo AC
computes with the unorientable quotient rather than committing.

\end{itemize}

The terrain vocabulary of Section~\ref{sec:terrain} then sorts the
traditions into three regimes of one theory.  A reduction order is
precisely a cost function whose terrain has no ridges, so greedy
substitution---which is what rewriting \emph{is}---cannot get stuck
on an interface tradeoff, and confluence supplies the unique bottom:
rewriting to normal form is local search that provably reaches the
canonical element.  At small width, greedy substitution fails on
interface tradeoffs, but a frontier-indexed dynamic program over a
few interfaces per goal is exact: the Selinger regime.  At large or
unknown width, equality
saturation~\cite{tate2009equality,willsey2021egg} defers
canonicalization entirely, keeping the whole class and extracting at
the end.  Equality saturation's own origin story---as the
non-destructive alternative to completion when orientation decisions
are premature---is the deferred-commitment move taken to its limit.
The multe-graph formalizes the middle regime, which is Selinger's:
defer exactly the distinctions the future can observe, and no more.
Neither rewriting tradition sits there---completion commits
everything and saturation defers everything---so what the
multe-graph offers rewriting is the middle discipline the database
tradition has practiced all along.

\begin{remark}[E-class analyses are labelings]
egg's e-class analyses are the functional presentation
(Remark~\ref{rem:presentations}) of multe-class labelings: an
e-node's analysis value is computed from its child classes' values,
a class's value is the join of its nodes' values, and the invariant
is maintained monotonically under merges---semantic factorization
plus join-semilattice structure, stated as a data-structure
invariant.  What the multe-graph adds is making the analysis
\emph{structural}: distinctions the analysis tracks become partition
boundaries that extraction can index, rather than side information
consulted by a cost function.  The familiar awkwardness of
property-sensitive extraction in egg---retrofitting frontiers onto a
one-best-term-per-class memo---is exactly the gap between carrying
labels and carrying interfaces.
\end{remark}

\paragraph{Completion as preprocessing.}

We believe the right way to operationalize the correspondence is to
run completion as a preprocessing pass.  Given a rewrite theory and
a reduction order compatible with the cost model, run ordered
completion.  It returns a convergent system $R$ together with an
unorientable residue $E$, and the two parts have different fates.

\begin{enumerate}

\item
\emph{Normalize with $R$ eagerly.}  Within width-one territory a
single representative is sufficient, so classes may be collapsed to
their $R$-normal forms before any multe-graph is built: no frontier,
no refinement hierarchy, no deferred commitment.

\item
\emph{Let $E$ localize the width.}  Every residue equation witnesses
a pair that the order cannot separate uniformly across contexts---a
candidate interface tradeoff.  Multe-class structure is needed only
there.

\item
\emph{Build the multe-graph over normal forms modulo $E$.}  The
interfaces to represent are the distinctions among $E$-equivalent
normal forms that remaining contexts can observe.

\end{enumerate}

On this architecture, completion is a width detector and localizer:
the convergent part of the theory is discharged by greedy rewriting,
while saturation and frontier-indexed extraction are reserved for
the residue, shrinking the e-graph before saturation begins.  This
also softens the hedge of Section~\ref{sec:width}: completion does
not choose a decomposition, but it improves the given one, rewriting
it into a form as monotone as the objective permits and confining
deferred commitment to the residue.  The
interleaving has a familiar feel: egg already alternates saturation
with congruence-based canonicalization (rebuilding); this adds
canonicalization by oriented rules to the same loop.

\begin{conjecture}[Completion measures width]
\label{conj:completion-width}
Over a ground-total reduction order, ordered completion succeeds
with an empty unorientable residue if and only if every goal has
width one under the order.  More generally, the distinctions that
require multe-class structure in the instantiated graph are exactly
those witnessed by residue equations, relative to the remaining
contexts of the instance.
\end{conjecture}

The conjecture would make completion failure a complexity statement
about the problem rather than a deficiency of the procedure---the
same reframing that Proposition~\ref{prop:trivial} performs for
Bellman sufficiency.

Finally, the critical-pair mechanism suggests the proof strategy for
Conjecture~\ref{conj:universality}.  Start from the coarsest
quotient, the full class; find \emph{critical contexts} that
distinguish members, the analogue of critical pairs; split the
class; iterate to a fixed point.  The fixed point should be the
continuation quotient.  This is Hopcroft-style partition refinement,
which is exactly how Myhill--Nerode is proved constructively---so
the completion connection is not merely an application of the
framework but a candidate algorithm for its central theorem.

\begin{openproblem}[Semantic completion]
\label{open:completion}
Define the completion procedure for semantic interfaces: critical
contexts as the analogue of critical pairs, class splitting as the
analogue of rule addition, termination and confluence of the
refinement process itself, and identification of its fixed point
with the continuation quotient.  Determine whether per-interface
normal forms---one canonical form per surviving multe-class, with a
critical-pair-style local coherence check across the refinement
hierarchy---coincide with an existing notion in rewriting modulo a
theory before coining a new one.
\end{openproblem}

\begin{remark}[Equality is not optimization]
One caveat keeps the correspondence honest.  Completion decides
equality; it is not natively an optimization procedure, and a normal
form is canonical, not necessarily the minimum of its class under
the (partial) reduction order.  The optimization reading tightens
when the order is ground-total, as superposition-based theorem
provers require.  The clean statement is therefore about
canonicalization coherence, with the optimization statement as a
corollary under ground-totality.
\end{remark}

%%%%%%%%%%%%%%%%%%%%%%%%%%%%%%%%%%%%%%%%%%%%%%%%%%%%%%%%%%%%%%%
\subsection{A Dictionary for Readers from Programming Languages}
\label{sec:types}
%%%%%%%%%%%%%%%%%%%%%%%%%%%%%%%%%%%%%%%%%%%%%%%%%%%%%%%%%%%%%%%

This closing subsection is a payoff of translation rather than of
new results.  In the outbound article it might be compressed into a
remark or a related-work discussion; we record it in full here
because it names several of our central notions with terms of art
that the programming-languages community has tested for
decades---and because, in one important case, the standard
vocabulary imports a warning we should heed.  The notation \(n:x\)
of Section~\ref{sec:goals-enodes} already anticipates the reading.

\begin{itemize}

\item
\emph{Goals are base types.}  The goal set \(X\) is a set of sorts;
a composition signature \(d=(x\Rightarrow x_1,\ldots,x_k)\) is a
constructor signature \(x_1\times\cdots\times x_k\to x\); e-nodes
are constructors; realizations are closed well-typed terms; and
\(\mathcal T(x)\) is the set of inhabitants of \(x\).  A concrete
semantic schema is a multi-sorted signature, and its solution space
is the term algebra.  Curry--Howard supplies the reading already
gestured at in Section~\ref{sec:goals-enodes}: a goal is a
proposition, realizations are its proofs, and extraction is
cost-aware proof search.

\item
\emph{Multe-classes are refinement
types}~\cite{freeman1991refinement}.  A multe-class is
\(\{\nu:x\mid\varphi(\nu)\}\), a refinement of the base type by a
semantic predicate, and the refinement hierarchy is a subtyping
lattice: \([R]_{\mathsf{sorted}(a)}<:[R]_{\mathsf{any}}\).  Child
references say that operators have refinement-typed signatures:
merge join demands
\([R]_{\mathsf{sorted}(a)}\times[S]_{\mathsf{sorted}(a)}\), while
hash join demands only the base types.  Interesting orders are a
refinement type system for physical plans.  Even the Cartesian-fiber
condition translates: an e-node whose concrete fiber requires
several rectangles is a constructor with an \emph{intersection
type}, one arrow per rectangle---so the size of the intersection is
the rectangle cover number of Remark~\ref{rem:rectangles}.

\item
\emph{Requirements are type ascriptions, and instantiation is
bidirectional type checking}~\cite{dunfield2021bidirectional}.  The
root requirement is an annotation at the root; liveness pushes
demanded types downward (checking mode); realizability synthesizes
provided types upward (synthesis mode); and the instantiated
multe-graph is what a bidirectional checker materializes.  This sits
beside the demand-transformation reading of
Section~\ref{sec:demand}: bidirectional typing and magic sets are
cousins.

\item
\emph{Continuation equivalence is contextual
equivalence}~\cite{morris1969}.
Definition~\ref{def:continuation} is Morris-style observational
equivalence: no program context can distinguish the terms.
Consequently, a Bellman-sufficient labeling is an \emph{adequate}
semantics, and an exact labeling is a \emph{fully abstract} one.

\end{itemize}

\begin{remark}[Full abstraction, and why we expect theorems]
The last identification restates the central conjectures in famous
words: Conjectures~\ref{conj:minimality}
and~\ref{conj:universality} assert that \emph{the fully abstract
model exists, is compositional, and is effectively presentable}.
The history of full abstraction for PCF~\cite{plotkin1977} supplies
the warning: the fully abstract model always exists
abstractly---quotient the term model, which is exactly
\(K_x\)---yet it can resist any syntax-independent, finitely
presented description for decades.  Our conjectures are precisely
the demand for a good presentation: a finite property vocabulary
with compositional transfer maps.  The reason to expect theorems
rather than a research program is that our setting is finite and
first-order---no higher-order functions, no divergence---which is
the regime in which full abstraction is tractable.  This observation
belongs in the outbound article: it preempts the obvious skeptical
question from readers who know the PCF story.
\end{remark}

The dictionary also imports proof technology.  The main obligation
of Conjecture~\ref{conj:universality} is that contextual equivalence
be a congruence---that composition descends to the quotients---and
the standard technique is Howe's method~\cite{howe1996proving}; in
our finite, first-order setting it should reduce to the partition
refinement already proposed as semantic completion
(Open Problem~\ref{open:completion}).  Finally, canonicalization
through a sufficient interface is recognizably normalization by
evaluation~\cite{berger1991inverse}: compute in the semantic model,
then read back the canonical form.  We take the density of these
correspondences as evidence that the framework is carving at known
joints.

%%%%%%%%%%%%%%%%%%%%%%%%%%%%%%%%%%%%%%%%%%%%%%%%%%%%%%%%%%%%%%%
\begin{thebibliography}{9}

\bibitem{selinger1979access}
P.~Griffiths Selinger, M.~M. Astrahan, D.~D. Chamberlin, R.~A. Lorie,
and T.~G. Price.
\newblock Access path selection in a relational database management
system.
\newblock In \emph{SIGMOD}, 1979.

\bibitem{heldkarp1962}
M.~Held and R.~M. Karp.
\newblock A dynamic programming approach to sequencing problems.
\newblock \emph{J. SIAM}, 10(1), 1962.

\bibitem{graefe1995cascades}
G.~Graefe.
\newblock The Cascades framework for query optimization.
\newblock \emph{IEEE Data Eng. Bull.}, 18(3), 1995.

\bibitem{willsey2021egg}
M.~Willsey, C.~Nandi, Y.~R. Wang, O.~Flatt, Z.~Tatlock, and
P.~Panchekha.
\newblock egg: Fast and extensible equality saturation.
\newblock \emph{Proc. ACM Program. Lang.}, 5(POPL), 2021.

\bibitem{bancilhon1986magic}
F.~Bancilhon, D.~Maier, Y.~Sagiv, and J.~D. Ullman.
\newblock Magic sets and other strange ways to implement logic
programs.
\newblock In \emph{PODS}, 1986.

\bibitem{tekle2011subsumptive}
K.~T. Tekle and Y.~A. Liu.
\newblock More efficient datalog queries: Subsumptive tabling beats
magic sets.
\newblock In \emph{SIGMOD}, 2011.

\bibitem{condie2008evita}
T.~Condie, D.~Chu, J.~M. Hellerstein, and P.~Maniatis.
\newblock Evita Raced: Metacompilation for declarative networks.
\newblock \emph{Proc. VLDB Endow.}, 1(1), 2008.

\bibitem{ioannidis1991left}
Y.~E. Ioannidis and Y.~C. Kang.
\newblock Left-deep vs. bushy trees: An analysis of strategy spaces
and its implications for query optimization.
\newblock In \emph{SIGMOD}, 1991.

\bibitem{dushnikmiller1941}
B.~Dushnik and E.~W. Miller.
\newblock Partially ordered sets.
\newblock \emph{American Journal of Mathematics}, 63(3), 1941.

\bibitem{trotter1992}
W.~T. Trotter.
\newblock \emph{Combinatorics and Partially Ordered Sets: Dimension
Theory}.
\newblock Johns Hopkins University Press, 1992.

\bibitem{yannakakis1982}
M.~Yannakakis.
\newblock The complexity of the partial order dimension problem.
\newblock \emph{SIAM J. Algebraic Discrete Methods}, 3(3), 1982.

\bibitem{johnson1988easy}
D.~S. Johnson, C.~H. Papadimitriou, and M.~Yannakakis.
\newblock How easy is local search?
\newblock \emph{J. Comput. Syst. Sci.}, 37(1), 1988.

\bibitem{papadimitriou1982combinatorial}
C.~H. Papadimitriou and K.~Steiglitz.
\newblock \emph{Combinatorial Optimization: Algorithms and
Complexity}.
\newblock Prentice-Hall, 1982.

\bibitem{knuthbendix1970}
D.~E. Knuth and P.~B. Bendix.
\newblock Simple word problems in universal algebras.
\newblock In \emph{Computational Problems in Abstract Algebra},
Pergamon Press, 1970.

\bibitem{newman1942}
M.~H.~A. Newman.
\newblock On theories with a combinatorial definition of
``equivalence''.
\newblock \emph{Annals of Mathematics}, 43(2), 1942.

\bibitem{bachmair1989completion}
L.~Bachmair, N.~Dershowitz, and D.~A. Plaisted.
\newblock Completion without failure.
\newblock In \emph{Resolution of Equations in Algebraic Structures},
volume~2, Academic Press, 1989.

\bibitem{baadernipkow1998}
F.~Baader and T.~Nipkow.
\newblock \emph{Term Rewriting and All That}.
\newblock Cambridge University Press, 1998.

\bibitem{tate2009equality}
R.~Tate, M.~Stepp, Z.~Tatlock, and S.~Lerner.
\newblock Equality saturation: A new approach to optimization.
\newblock In \emph{POPL}, 2009.

\bibitem{morris1969}
J.~H. Morris.
\newblock \emph{Lambda-Calculus Models of Programming Languages}.
\newblock PhD thesis, MIT, 1969.

\bibitem{plotkin1977}
G.~D. Plotkin.
\newblock LCF considered as a programming language.
\newblock \emph{Theoretical Computer Science}, 5(3), 1977.

\bibitem{freeman1991refinement}
T.~Freeman and F.~Pfenning.
\newblock Refinement types for ML.
\newblock In \emph{PLDI}, 1991.

\bibitem{dunfield2021bidirectional}
J.~Dunfield and N.~Krishnaswami.
\newblock Bidirectional typing.
\newblock \emph{ACM Computing Surveys}, 54(5), 2021.

\bibitem{howe1996proving}
D.~J. Howe.
\newblock Proving congruence of bisimulation in functional
programming languages.
\newblock \emph{Information and Computation}, 124(2), 1996.

\bibitem{berger1991inverse}
U.~Berger and H.~Schwichtenberg.
\newblock An inverse of the evaluation functional for typed
$\lambda$-calculus.
\newblock In \emph{LICS}, 1991.

\end{thebibliography}

\end{document}
